%%%%%%%%%%%%%%%%%%%%%%%%%%%%%%%%%%%%%%%%%%%%%%%%%%%%%%%%%%%%%%%%
%
% Apuntes de la asignatura Análisis Matemático I.
% Doble Grado de Informática y Matemáticas.
% Universidad de Granada.
% Curso 2016/17.
% 
% 
% Colaboradores:
% Javier Sáez (@fjsaezm)
% Daniel Pozo (@danipozodg)
% Pedro Bonilla (@pedrobn23)
% Guillermo Galindo
% Antonio Coín (@antcc)
% Sofía Almeida (@SofiaAlmeida)
%
% Agradecimientos:
% Andrés Herrera (@andreshp) y Mario Román (@M42) por
% las plantillas base.
%
% Sitio original:
% https://github.com/libreim/apuntesDGIIM/
%
% Licencia:
% CC BY-NC-SA 4.0 (https://creativecommons.org/licenses/by-nc-sa/4.0/)
%
%%%%%%%%%%%%%%%%%%%%%%%%%%%%%%%%%%%%%%%%%%%%%%%%%%%%%%%%%%%%%%%


%------------------------------------------------------------------------------
%   ACKNOWLEDGMENTS
%------------------------------------------------------------------------------

%%%%%%%%%%%%%%%%%%%%%%%%%%%%%%%%%%%%%%%%%%%%%%%%%%%%%%%%%%%%%%%%%%%%%%%%
% Plantilla básica de Latex en Español.
%
% Autor: Andrés Herrera Poyatos (https://github.com/andreshp) 
%
% Es una plantilla básica para redactar documentos. Utiliza el paquete  fancyhdr para darle un
% estilo moderno pero serio.
%
% La plantilla se encuentra adaptada al español.
%
%%%%%%%%%%%%%%%%%%%%%%%%%%%%%%%%%%%%%%%%%%%%%%%%%%%%%%%%%%%%%%%%%%%%%%%%%

%%%
% Plantilla de Trabajo
% Modificación de una plantilla de Latex de Frits Wenneker para adaptarla 
% al castellano y a las necesidades de escribir informática y matemáticas.
%
% Editada por: Mario Román
%
% License:
% CC BY-NC-SA 3.0 (http://creativecommons.org/licenses/by-nc-sa/3.0/)
%%%

%%%%%%%%%%%%%%%%%%%%%%%%%%%%%%%%%%%%%%%%
% Short Sectioned Assignment
% LaTeX Template
% Version 1.0 (5/5/12)
%
% This template has been downloaded from:
% http://www.LaTeXTemplates.com
%
% Original author:
% Frits Wenneker (http://www.howtotex.com)
%
% License:
% CC BY-NC-SA 3.0 (http://creativecommons.org/licenses/by-nc-sa/3.0/)
%
%%%%%%%%%%%%%%%%%%%%%%%%%%%%%%%%%%%%%%%%%


% Tipo de documento y opciones.
\documentclass[11pt, a4paper, titlepage]{article}


%---------------------------------------------------------------------------
%   PAQUETES
%---------------------------------------------------------------------------

% Idioma y codificación para Español.
\usepackage[utf8]{inputenc}
\usepackage[spanish, es-tabla, es-lcroman, es-noquoting]{babel}
\selectlanguage{spanish} 
%\usepackage[T1]{fontenc}

% Fuente utilizada.
\usepackage{courier}    % Fuente Courier.
\usepackage{microtype}  % Mejora la letra final de cara al lector.

% Diseño de página.
\usepackage{fancyhdr}   % Utilizado para hacer títulos propios.
\usepackage{lastpage}   % Referencia a la última página.
\usepackage{extramarks} % Marcas extras. Utilizado en pie de página y cabecera.
\usepackage[parfill]{parskip}    % Crea una nueva línea entre párrafos.
\usepackage{geometry}            % Geometría de las páginas.

% Símbolos y matemáticas.
\usepackage{amssymb, amsmath, amsthm, amsfonts, amscd}
\usepackage{upgreek}

% Otros.
\usepackage{enumitem}   % Listas mejoradas.
\usepackage[hidelinks]{hyperref}
\usepackage{graphicx}            % Gráficos.

%---------------------------------------------------------------------------
%   OPCIONES PERSONALIZADAS
%---------------------------------------------------------------------------

% Redefinir letra griega épsilon.
\let\epsilon\upvarepsilon

% Formato de texto.
\linespread{1.1}            % Espaciado entre líneas.
\setlength\parindent{0pt}   % No indentar el texto por defecto.
\setlist{leftmargin=.5in}   % Indentación para las listas.

% Estilo de página.
\pagestyle{fancy}
\fancyhf{}
\geometry{left=3cm,right=3cm,top=3cm,bottom=3cm,headheight=1cm,headsep=0.5cm}   % Márgenes y cabecera.

% Ruta donde buscar gráficos
\graphicspath{{../Recursos/Plantillas/}, {Recursos/Plantillas/}}

% Redefinir entorno de demostración (reducir espacio superior)
\makeatletter
\renewenvironment{proof}[1][\proofname] {\vspace{-15pt}\par\pushQED{\qed}\normalfont\topsep6\p@\@plus6\p@\relax\trivlist\item[\hskip\labelsep\it#1\@addpunct{.}]\ignorespaces}{\popQED\endtrivlist\@endpefalse}
\makeatother

% Aumentar el tamaño del interlineado
\linespread{1.3}
%---------------------------------------------------------------------------
%   COMANDOS PERSONALIZADOS
%---------------------------------------------------------------------------

% Valor absoluto: \abs{}
\providecommand{\abs}[1]{\lvert#1\rvert}    

% Fracción grande: \ddfrac{}{}
\newcommand\ddfrac[2]{\frac{\displaystyle #1}{\displaystyle #2}}

% Texto en negrita en modo matemática: \bm{}
\newcommand{\bm}[1]{\boldsymbol{#1}}

% Línea horizontal.
\newcommand{\horrule}[1]{\rule{\linewidth}{#1}}

% Renombrar la R de los reales

\newcommand{\R}{\mathbb{R}}

%---------------------------------------------------------------------------
%   CABECERA Y PIE DE PÁGINA
%---------------------------------------------------------------------------

% Cabecera del documento.
\renewcommand\headrule{
	\begin{minipage}{1\textwidth}
		\hrule width \hsize 
	\end{minipage}
}

% Texto de la cabecera.
\lhead{\subject}  % Izquierda.
\chead{}            % Centro.
\rhead{\docauthor}    % Derecha.

% Pie de página del documento.
\renewcommand\footrule{                                 
	\begin{minipage}{1\textwidth}
		\hrule width \hsize   
	\end{minipage}\par
}

% Texto del pie de página.
\lfoot{}                                                 % Izquierda
\cfoot{}                                                 % Centro.
\rfoot{Página\ \thepage\ de\ \protect\pageref{LastPage}} % Derecha.


%---------------------------------------------------------------------------
%   ENTORNOS PARA MATEMÁTICAS
%---------------------------------------------------------------------------

% Nuevo estilo para definiciones.
\newtheoremstyle{definition-style} % Nombre del estilo.
{10pt}               % Espacio por encima.
{10pt}               % Espacio por debajo.
{}                   % Fuente del cuerpo.
{}                   % Identación.
{\bf}                % Fuente para la cabecera.
{.}                  % Puntuación tras la cabecera.
{.5em}               % Espacio tras la cabecera.
{\thmname{#1}\thmnumber{ #2}\thmnote{ (#3)}}     % Especificación de la cabecera (actual: nombre en negrita).

% Nuevo estilo para notas.
\newtheoremstyle{remark-style} 
{10pt}                
{10pt}                
{}                   
{}                   
{\itshape}          
{.}                  
{.5em}               
{}                  

% Nuevo estilo para teoremas y proposiciones.
\newtheoremstyle{theorem-style}
{10pt}                
{10pt}                
{\itshape}           
{}                  
{\bf}             
{.}                
{.5em}               
{\thmname{#1}\thmnumber{ #2}\thmnote{ (#3)}}                   

% Nuevo estilo para ejemplos.
\newtheoremstyle{example-style}
{10pt}                
{10pt}                
{}                  
{}                   
{\scshape}              
{:}                 
{.5em}               
{}                   

% Teoremas, proposiciones y corolarios.
\theoremstyle{theorem-style}
\newtheorem*{nth}{Teorema}
\newtheorem*{nprop}{Proposición}
\newtheorem{ncor}{Corolario}

% Definiciones.
\theoremstyle{definition-style}
\newtheorem*{ndef}{Definición}

% Notas.
\theoremstyle{remark-style}
\newtheorem*{nota}{Nota}

% Ejemplos.
\theoremstyle{example-style}
\newtheorem*{ejemplo}{Ejemplo}

% Listas ordenadas con números romanos (i), (ii), etc.
\newenvironment{nlist}
{\begin{enumerate}
\renewcommand\labelenumi{(\emph{\roman{enumi})}}}
{\end{enumerate}}

% División por casos con llave a la derecha.
\newenvironment{rcases}
  {\left.\begin{aligned}}
  {\end{aligned}\right\rbrace}

%---------------------------------------------------------------------------
%   PÁGINA DE TÍTULO
%---------------------------------------------------------------------------

% Título del documento.
\newcommand{\subject}{Análisis Matemático I}

% Autor del documento.
\newcommand{\docauthor}{Doble Grado de Informática y Matemáticas}

% Título
\title{
  \normalfont \normalsize 
  \textsc{Universidad de Granada} \\ [25pt]    % Texto por encima.
  \horrule{0.5pt} \\[0.4cm] % Línea horizontal fina.
  \huge \subject\\ % Título.
  \horrule{2pt} \\[0.5cm] % Línea horizontal gruesa.
}

% Autor.
\author{\Large{\docauthor}}

% Fecha.
\date{\vspace{-1.5em} \normalsize Curso 2016/17}


%---------------------------------------------------------------------------
%   COMIENZO DEL DOCUMENTO
%---------------------------------------------------------------------------
\begin{document}

\maketitle  % Título.
\tableofcontents    % Índice
\vfill
\begin{center}
\includegraphics{by-nc-sa}  % Licencia.
\end{center}
\newpage

%----------------------------
%   Introducción.
%----------------------------
\section*{Introducción.}

El objetivo de este curso es el estudio de las funciones de varias variables, es decir, de funciones $f: \mathbb{R}^N \longrightarrow \mathbb{R}^M$. Para ello, empezaremos caracterizando el espacio $\mathbb{R}^N$, y proseguiremos intentando traspasar los resultados principales sobre funciones reales de variable real a nuestro campo de estudio, así como enunciando otros nuevos.

Es por esto que es fundamental haber cursado con aprovechamiento las asignaturas de \emph{Cálculo I y II}, que tratan exclusivamente sobre funciones reales de variable real.

Aunque nos centraremos en funciones en el espacio $\mathbb{R}^N$, muchos de los resultados que obtendremos son igual de válidos en un espacio métrico en general, e incluso en espacios topológicos.

\newpage


%--------------------------------------
%   Topología de un espacio métrico.
%--------------------------------------
\section{Topología de un espacio métrico.}

%%% El espacio métrico R^N.
\subsection{Concepto de espacio métrico. El espacio métrico $\mathbb{R}^N$.}

\begin{ndef}[Espacio métrico]
Consideremos un conjunto $X$ cualquiera, y una aplicación \mbox{$d:X\times X \longrightarrow \mathbb{R}$} que cumple las siguientes propiedades:

\begin{nlist}
\item $d(x,y) \ge 0\ \ \forall x,y \in X$.

\item $d(x,y) = 0 \iff x = y\ \ \forall x,y \in X$.

\item $d(x,y) = d(y,x)\ \ \forall x,y \in X$.

\item $d(x,y) \leq d(x,z) + d(z,y)\ \ \forall x,y,z \in X. \quad \quad(desigualdad\ triangular)$

\end{nlist}

Entonces, se dice que el par $(X,d)$ es un \emph{espacio métrico}.
\end{ndef}



\begin{nota}
En adelante, entenderemos $\mathbb{R}^N$ como el espacio métrico $(\mathbb{R}^N,d)$, siendo $d$ la distancia usual \textbf{(distancia euclídea)} dada por: $$d(x,y) = \sqrt{\sum_{i=1}^N (y_i - x_i)^2}\ \ \forall x,y\in \mathbb{R}^N.$$

Existen otras distancias en $\mathbb{R}^N$. Las más destacadas son las siguientes:

\begin{nlist}
\item $\displaystyle d_1(x,y) = \sum_{i=1}^N \abs{x_i - y_i}\ \ \forall x,y \in \mathbb{R}^N$.

\item $\displaystyle d_{\infty}(x,y) = m\acute{a}x \{\abs{x_i - y_i}: \ i=1,\dots,N \}\ \ \forall x,y\in \mathbb{R}^N$.

\item $\displaystyle d_p(x,y) = \left( \sum_{i=1}^N \abs{x_i - y_i}^p \right)^{1/p}\ \ \forall x,y\in \mathbb{R}^N$.\\
\end{nlist}

\end{nota}



\begin{ndef}
Sean $(X,d)$ y $(X,d')$ dos espacios métricos sobre un mismo conjunto $X$. Se dice que las distancias $d$ y $d'$ son \textit{equivalentes} si, y solo si, $$\exists k_1,k_2 > 0 :\ k_1d(x,y)\le d'(x,y) \le k_2d(x,y)\ \ \forall x,y\in X.$$
\end{ndef}



\begin{nprop}
En $\mathbb{R}^N$, todas las distancias mencionadas anteriormente son equivalentes entre sí. En particular, la distancia euclídea es equivalente a todas ellas.
\end{nprop}


%%% Conceptos topológicos.
\subsection{Conceptos topológicos.} 

\begin{ndef}[Bola abierta]
Sea $(X,d)$ un espacio métrico, y fijemos un $x\in X$ y un $\epsilon > 0$. Se llama \emph{bola abierta de centro $x$ y radio $\epsilon$} al conjunto $B(x,\epsilon) = \{ y\in X \ | \ d(x,y)<\epsilon\}$.
\end{ndef}



\begin{ndef}[Bola cerrada]
De forma análoga, se define la \emph{bola cerrada de centro $x$ y radio $\epsilon$} como el conjunto $\overline{B}(x,\epsilon) = \{y\in X \ | \ d(x,y)\leq \epsilon \}$.
\end{ndef}



\begin{ndef}[Conjunto abierto]
Sea $(X,d)$ un espacio métrico, y sea $A\subseteq X$. Decimos que \mbox{$A\ es\ abierto \iff \forall a \in A\ \exists \epsilon > 0: B(x,\epsilon) \subseteq A$}.	
\end{ndef}



\begin{nprop}
Sea $(X,d)$ un espacio métrico. Entonces, $\forall x \in X \ \forall \epsilon > 0$ se tiene que $B(x,\epsilon)$ es un conjunto abierto.
\end{nprop}

\begin{proof}

Sea $x \in B(x_0,\epsilon_0) $ arbitrario. Para demostrar que $B(x_0,\epsilon_0)$ es un abierto, tenemos que encontrar un $\epsilon > 0$ tal que $B(x,\epsilon) \subseteq B(x_0,\epsilon_0)  $, y por lo tanto comprobar que se verifica que $\forall y \in B(x,\epsilon) \Rightarrow y \in B(x_0,\epsilon_0).$

Sea $y \in B(x, \epsilon )$ cualquiera. Consideremos $r = d(x, x_0)$, y tomemos $\epsilon = \epsilon_0 - r$. Queremos demostrar que $y \in B(x_0, \epsilon_0)$. Para ello, veamos que $d(x_0, y) < \epsilon_0$. En efecto, por la desigualdad triangular se cumple que: $$d(x_0, y) \leq d(x,x_0) + d(x,y) < r + \epsilon = r + \epsilon_0 - r = \epsilon_0 $$
Luego queda demostrado que $y \in B(x_0, \epsilon_0)$, y por tanto podemos afirmar que para todo punto $x \in B(x_0, \epsilon_0)$ se puede encontrar una bola abierta centrada en él, tal que todos sus puntos están en el conjunto de origen.
\end{proof}



\begin{nprop}
Sea $(X,d)$ un espacio métrico. Entonces, se verifican las siguientes propiedades:

\begin{nlist}
\item $Si\ \{A_\lambda \ | \ \lambda \in \Lambda \}$ es una familia de subconjuntos abiertos de $X$, entonces $\displaystyle \bigcup_{\lambda \in \Lambda} A_\lambda$ es abierto.

\item Si $\{A_1,\dots, A_n\}$ es una familia finita de abiertos de $X$, entonces $\displaystyle \bigcap_{i=1}^n A_i$ es abierto.

\item $X,\emptyset$ son abiertos.
\end{nlist}

\end{nprop}



\begin{ndef}[Punto interior]
Sea $(X,d)$ un espacio métrico, y consideremos $A\subseteq X$, $a\in A$. Se dice que $a$ \emph{es un punto interior de} $A$ si, y solo si, $\exists \epsilon_0 > 0: B(a,\epsilon_0)\subseteq A$.  Definimos $int(A) = \mathring{A} = \{ a\in A \ | \ a\ es\ punto\ interior\ de\ A\}$.
\end{ndef}



\begin{nprop}
Sea $(X,d)$ un espacio métrico, y $A\subseteq X$. Entonces, se verifican las siguientes propiedades:

\begin{nlist}
\item $\mathring{A} \subseteq A$.

\item $\mathring{A}$ es abierto.

\item Si $B\subseteq A$ es un subconjunto abierto de $A$, entonces $B \subseteq \mathring{A}$. Es decir, $\mathring{A}$ es el abierto más grande contenido en $A$.

\item $\displaystyle \mathring{A}  = \bigcup \{ B\subseteq A \ | \ B\ es\ abierto \}$.

\item $A$ es abierto $\iff \mathring{A} =A$.

\item $int(int(A)) = int(A).$

\item Si $A\subseteq B$, entonces $\mathring{A} \subseteq \mathring{B}$.
\end{nlist}

\end{nprop}



\begin{ndef}[Conjunto cerrado]
Sea $(X,d)$ un espacio métrico, y $F\subseteq X$. Se dice que el conjunto $F\ es\ cerrado \iff X-F\ es\ abierto$.
\end{ndef}



\begin{nprop}
Sea $(X,d)$ un espacio métrico. Entonces, $\forall x\in X \ \forall \epsilon > 0$ se tiene que $\overline{B}(x,\epsilon)$ es un conjunto cerrado.
\end{nprop}



\begin{nprop}
Sea $(X,d)$ un espacio métrico. Entonces, se verifican las siguientes propiedades:

\begin{nlist}
\item $Si\ \{F_\lambda \ | \ \lambda \in \Lambda \}$ es una familia de cerrados de $X$, entonces $\displaystyle \bigcap_{\lambda \in \Lambda} F_\lambda$ es cerrado.

\item Si $\{F_1,\dots, F_n\}$ es una familia finita de cerrados de $X$, entonces $\displaystyle \bigcup_{i=1}^n F_i$ es cerrado.

\item $X,\emptyset$ son cerrados.
\end{nlist}

\end{nprop}



\begin{ndef}[Clausura]
Sea $(X,d)$ un espacio métrico. Se llama \textit{clausura o cierre de A} al conjunto $\overline{A} = X - int(X-A)$.
\end{ndef}



\begin{nprop}
Sea $(X,d)$ un espacio métrico, y $A\subseteq X$. Entonces, se verifican las siguientes propiedades:
\begin{nlist}
\item $A \subseteq \overline{A}$.

\item $\overline{A}$ es cerrado.

\item Si $B\subseteq X$ es un subconjunto cerrado de $X$ tal que $A\subseteq B$, entonces $\overline{A} \subseteq B$. Es decir, $\overline{A}$ es el cerrado más pequeño que contiene a $A$.

\item $\displaystyle \overline{A}  = \bigcap \{ F\subseteq X \ | \ F\ es\ cerrado\ y\ A\subseteq F \}$.

\item $A$ es cerrado $\iff \overline{A} = A$.

\item $\overline{\overline{A}} = \overline{A}.$

\item Si $A\subseteq B$, entonces $\overline{A} \subseteq \overline{B}$.
\end{nlist}

\end{nprop}



\begin{ndef}[Frontera]
Sea $(X,d)$ un espacio métrico, y $A\subseteq X$. Llamamos \textit{frontera de A} al conjunto $\partial A = \overline{A}-\mathring{A}$.	
\end{ndef}



\begin{nprop}
Sea $(X,d)$ un espacio métrico, y $A\subseteq X$. Entonces, se verifica lo siguiente:
$x\in \partial A \iff \forall \epsilon > 0\ B(x,\epsilon)\cap A \neq \emptyset \ y\ B(x,\epsilon)\cap (X-A) \neq \emptyset$.
\end{nprop}



\begin{ndef}[Punto de acumulación]
Sea $(X,d)$ un espacio métrico, y $A\subseteq X$. Dado $x\in X$, decimos que \textit{x es punto de acumulación de} $A \iff \forall \epsilon > 0\ B(x,\epsilon)\cap (A-\{x\})\neq \emptyset$. Definimos $A' = \{ x\in X \ | \ x\ es\ punto\ de\ acumulaci\acute{o}n\ de\ A \}$.
\end{ndef}



\begin{nprop}
Sea $(X,d)$ un espacio métrico. Entonces, se verifican las siguientes afirmaciones:

\begin{nlist}
\item $\mathring{A} = X - \overline{X-A}$

\item $\overline{A} = A \cup \partial A$.

\item $\overline{A} = A \cup A'$

\item $X = int(A) \cup \partial A \cup int(X-A)$. Además, la unión es disjunta dos a dos.
\end{nlist}

\end{nprop}

\newpage


%--------------------------------------
%   Sucesiones en R^N.
%--------------------------------------
\section{Sucesiones en $\mathbb{R}^N$.}



\begin{ndef}[Sucesión en $\mathbb{R}^N$]
Una sucesión en $\mathbb{R}^N$ es una aplicación $x: \mathbb{N} \longrightarrow \mathbb{R}^N$ que a cada $n\in \mathbb{N}$ le hace corresponder un $x(n) \in \mathbb{R}^N$. Por simplicidad, al elemento imagen de $n$ se le denomina $x_n$, y la aplicación $x$ se denota $\{x_n\}$.
\end{ndef}



\begin{ndef}[Convergencia de sucesiones]
Sea $(X,d)$ un espacio métrico, $A\subseteq X$ y $x\in X$. Decimos que una sucesión $\{x_n\}$ de puntos de $A$ converge a $x$ si, y solo si: $$ \forall \epsilon > 0\ \ \exists n_o \in \mathbb{N}: \ n\ge n_o \Rightarrow d(x_n,x) < \epsilon.$$
\end{ndef}



\begin{nota}
Este concepto no depende de la distancia equivalente elegida.
\end{nota}



\begin{nprop}
Sea $A\subseteq \mathbb{R}^N$, $x\in \mathbb{R}^N$, y $\{x_n\}$ una sucesión de puntos de $A$. Adoptemos la notación $x_n = (x_n^1, x_n^2,\dots,x_n^N)$, y $x=(x^1, x^2,\dots, x^N)$. Entonces, se verifica que:
$$\{x_n\} \rightarrow x \iff \{x_n^j\} \rightarrow x^j.$$
\end{nprop}



%\begin{proof}
%\fbox{$\impliedby$}

%Para esta demostración vamos a usar la distancia del máximo en lugar de la usual (todas las distancias son equivalentes en $\mathbb{R}^N$).\\

%Si $\forall i=1, ..., N$ tenemos que $\forall \varepsilon_i > 0\quad \exists m_i : n \ge m_i \implies |x^i_n - x^i| < \varepsilon_i$, se sigue que \[\displaystyle\max_{i=1,...,N} |x^i_n - x^i| < \displaystyle\max_{i=1,...,N} \varepsilon_i\]
%luego $d_\inf



\begin{ndef}
Sea $(X,d)$ un espacio métrico, y $x\in X$. Consideremos, para cada $n\in \mathbb{N}$, un punto $a_n \in X$. Entonces, decimos que $d(a_n,x) \rightarrow 0$ $\iff \{a_n\} \rightarrow x$.
\end{ndef}



\begin{ndef}[Conjunto acotado]
Sea $A \subseteq \mathbb{R}^N$. Decimos que $A$ \textit{está acotado} si, y solo si, $\exists R>0: A\subseteq B(0,R)$.
\end{ndef}



\begin{ndef}[Sucesión acotada]
Sea $\{x_n\}$ una sucesión de puntos de $R^N$. Entonces, decimos que $\{x_n\}$ \textit{está acotada} sí, y solo sí, $\{x_n\ | \ n\in \mathbb{N}\}$ está acotado.
\end{ndef}



\begin{nprop}
Si una sucesión $\{x_n\} \subseteq \mathbb{R}^N$ es acotada, entonces $\forall i=1,\dots,n$ la sucesión $\{x_n^i\}$ es acotada (en $\mathbb{R}$).
\end{nprop}



\begin{nota}
Si un conjunto $A\subseteq \mathbb{R}^N$ es acotado, entonces cualquier sucesión de puntos de $A$ es acotada.
\end{nota}



\begin{nth}[Bolzano-Weierstrass]
Sea $\{x_n\}\subseteq \mathbb{R}^N$ acotada. Entonces, existe una sucesión parcial suya $\{x_{\sigma_{(n)}}\}$ convergente.
\end{nth}

\begin{proof}
	Notaremos $x_n = (x_n^1, \dots, x_n^N)$. Como $\{x_n^1\}$ es acotada en $\mathbb{R}$, existe $\sigma_1 : \mathbb{N} \rightarrow \mathbb{N}$ estrictamente creciente tal que $\{x_{\sigma_1(n)}^1\}$ es convergente.

Ahora, como $\{x_n^2\}$ es acotada, $\{x_{\sigma_1(n)}^2\}$ también es acotada, y existe $\sigma_2 : \mathbb{N} \rightarrow \mathbb{N}$ estrictamente creciente tal que $\{x_{(\sigma_2\circ\sigma_1)(n)}^1\}$ es convergente.

Procediendo de esta forma con cada componente de $x_n$, obtenemos $\sigma_1, \dots, \sigma_N$, y\\ $\{x_{\sigma_1(n)}^1\}, \{x_{(\sigma_2\circ\sigma_1)(n)}^2\}, \dots, \{x_{(\sigma_N\circ\dots\circ\sigma_2\circ\sigma_1)(n)}^N\}$ sucesiones convergentes en $\mathbb{R}$. Al ser $\sigma_i$ estrictamente creciente $\forall i=1,\dots,N$, $\{x_{(\sigma_N(n)\circ\dots\circ\sigma_{i+1}\sigma_i\circ\dots\circ\sigma_1)(n)}^i\}$ también es convergente (toda sucesión parcial de una sucesión convergente es convergente).

Así, tomando $\sigma = \sigma_1\circ\dots\circ\sigma_N$, $\{x_{\sigma(n)}\}$ es convergente.
\end{proof}


\begin{ndef}[Sucesión de Cauchy]
Sea $\{x_n\}\subseteq \mathbb{R}^N$. Decimos que $\{x_n\}$ es una \textit{sucesión de Cauchy} si $$ \forall \epsilon > 0\ \ \exists n_o \in \mathbb{N}: \ n,m\ge n_o \Rightarrow d(x_n,x_m) < \epsilon$$
\end{ndef}



\begin{nth} [$\bm{\mathbb{R}^N}$ es completo]
Sea $\{x_n\}\subseteq \mathbb{R}^N$. Entonces: $$\{x_n\}\ es\ de\ Cauchy\ \iff \{x_n\}\ es\ convergente.$$
\end{nth}



\begin{nprop}
Sea $\{x_n\} \subseteq \mathbb{R}^N\ con\ \{x_n\} \rightarrow x \in \mathbb{R}^N$. Entonces, toda sucesión parcial de $\{x_n\}$ es convergente a $x$.
\end{nprop}
\begin{proof}
	Sea $\sigma :\mathbb N \to \mathbb N$ una aplicación tal que $\sigma(n+1) > \sigma(n) > n \quad \forall n \in \mathbb N$. Ahora, $\{x_{\sigma(n)}\}$ es una sucesión parcial de $\{x_n\}$. Como $x_n$ es convergente:
	\[
	\forall \epsilon > 0 \ \ \ \exists n_0 \in \mathbb N : \ \ n \geq n_0 \implies d(x_n,x) < \epsilon
	\]
	Pero $\sigma(n) \geq n \ \ \forall n$ así que
	\[
	\forall \epsilon > 0 \ \ \ \exists n_0 \in \mathbb N : \ \ \sigma(n) \geq n \geq n_0 \implies d(x_{\sigma(n)},x) < \epsilon
	\]
	Luego $\{x_{\sigma(n)}\}\to x$
\end{proof}

\newpage


%-----------------------------------
%   Funciones continuas en R^N.
%-----------------------------------
\section{Funciones continuas en $\mathbb{R}^N$.}



\begin{ndef}[Función continua]
Sea $\emptyset \ne A\subseteq \mathbb{R}^N$, $f: A \longrightarrow \mathbb{R}^M$ y $a \in A$. Decimos que \textit{f es continua en $a$} si, y solo si: $$\forall \epsilon > 0\ \ \exists \delta > 0: \ x\in A, \ d(x,a)<\delta \Rightarrow d(f(x),f(a))<\epsilon.$$
Además, se dice que $f$ es continua si lo es en todos sus puntos.	
\end{ndef}



\begin{nprop}[Caracterización de continuidad]
Sea $\emptyset \ne A \subseteq \mathbb{R}^N$, y $f:A\longrightarrow \mathbb{R}^M$. Entonces: $$f\ es\ continua\ en\ a \iff \forall \{x_n\}\subseteq A\ con\ \{x_n\} \rightarrow a \Rightarrow \{f(x_n)\} \rightarrow f(a).$$
\end{nprop}



\begin{ndef}[Continuidad uniforme]
Sea $\emptyset \ne A \subseteq \mathbb{R}^N$, $f:A \longrightarrow \mathbb{R}^M$. Se dice que $f$ es uniformemente continua si, y solo si: $$\forall \epsilon > 0 \ \ \exists \delta > 0 : \ x,y \in A,\ d(x,y) < \delta \Rightarrow d(f(x),f(a)) < \epsilon.$$
\end{ndef}


\begin{ndef}[Conjunto compacto]
Sea $(X,d)$ un espacio métrico, y sea $\emptyset \ne A \subseteq X$. $$A\ es\ compacto \iff \forall \{x_n\} \subseteq A\ \ \exists \{x_{\sigma(n)}\} \rightarrow x\in A.$$
\end{ndef}

\begin{ndef}[Recubrimiento abierto]
	Sea $A \subseteq \mathbb{R}^N$. Se dice que una familia $\{O_i, i\in I\}$ de abiertos es un \emph{recubrimiento abierto} de $A$ si
	\[
		A \subseteq \bigcup_{i\in I} O_i
	\]
	
	También, si $R_1$ y $R_2$ son recubrimientos abiertos de $A$ y $R_1 \subseteq R_2$, se dice que $R_1$ es un \emph{subrecubrimiento abierto} de $R_2$.
\end{ndef}

\begin{nprop}
	Sea $A \subseteq \mathbb{R}^N$ y $\bigcup_{i\in I} O_i$ un recubrimiento cualquiera de $A$. Entonces, $A$ es compacto si $\exists n\in \mathbb N$ tal que 
	\[
	A \subseteq \bigcup_{i=1}^n O_i
	\]
	es decir, existe un subrecubrimiento de $A$ finito.
\end{nprop}

\begin{nprop}[Caracterización de cerrados]
Sea $(X,d)$ un espacio métrico, y $A\subseteq X$. Entonces, son equivalentes:

\begin{nlist}
\item $A$ es cerrado.
\item $\forall \{x_n\} \subseteq A$ convergente a un $x \in X$, se verifica que $x\in A$.
\end{nlist}

\end{nprop}

\begin{proof} 
Veamos las dos implicaciones:

$\displaystyle \boxed{\Rightarrow}\ $ Supongamos $A \subseteq X$ un conjunto cerrado. Entonces, $X - A$ es abierto. Sea $\{x_n\}$ una sucesión de puntos de $A$ que converge a un $x\in X$. Para comprobar que, de hecho, $x\in A$, argumentamos por reducción al absurdo:

\underline{Supongamos $x\notin A$}. Entonces, $x\in X - A$, y por ser este último conjunto abierto, encontramos un $\epsilon > 0\ tal\ que\ B(x,\epsilon)\subseteq (X-A)$. Pero por ser $x$ el límite de la sucesión $\{x_n\}$, se tiene que $\exists n_o \in \mathbb{N}: n\geq n_o \Rightarrow d(x_n,x)<\epsilon$. Es decir, a partir de cierto índice en adelante, $x_n \in B(x,\epsilon)\ con\ x_n \in A\ \forall n \in \mathbb{N}$. Esto se contradice con el hecho de que $B(x,\epsilon)\subseteq (X-A)$, pues encontramos en dicha bola puntos $x_n$ que no pertenecen a $X-A$.

Por tanto, concluimos que $x\in A$.

$\displaystyle \boxed{\Leftarrow}\ $ Sea $A\subseteq X$, y supongamos que se verifica que $\forall \{x_n\} \subseteq A\ tal\ que\ \{x_n\} \rightarrow x \in X,\ se\ tiene\ que\ x\in A$. Para ver que $A$ es cerrado, utilizaremos la siguiente caracterización de conjuntos cerrados:
\vspace{0.5em}
$$A\ es\ cerrado\ \iff \overline{A} = A $$

Si recordamos, se define la frontera de $A$ como $\partial A = \overline{A} - \mathring{A}$. Por tanto, la equivalencia anterior quedaría así: $A\ es\ cerrado \iff \partial A \cup \mathring{A} = A$. Para comprobar esta última igualdad, veamos las dos inclusiones:

\begin{description}
\item $\displaystyle \boxed{\subseteq}\ $ Sabemos por la definición del conjunto de puntos interiores de $A$, que $\mathring{A} \subseteq A$. \\ Comprobemos entonces que $\partial A \subseteq A$:

Sea $x\in \partial A$ cualquiera. Por una caracterización de la frontera de $A$, sabemos que $\forall \epsilon > 0\ B(x,\epsilon)\cap A \neq \emptyset$. Si tomamos $\epsilon = \frac{1}{n} > 0\ con \ n\in \mathbb{N}$, tenemos que $B(x,\frac{1}{n})\cap A \neq \emptyset$, es decir, $\exists a_n \in B(x,\frac{1}{n})\cap A\ tal\ que\ d(x,a_n)<\epsilon = \frac{1}{n}$. Podemos construir entonces, para cada $n \in \mathbb{N}$, la sucesión $\{a_n\}$.

Así, se tiene que $0 < d(x,a_n) < \frac{1}{n}\ \forall n\in \mathbb{N}$, de donde concluimos que $d(x,a_n) \rightarrow 0$. Por definición, esto significa que $\{a_n\} \rightarrow x$, lo que por hipótesis implica, al ser $\{a_n\}$ una sucesión convergente de puntos de $A$, que  $x\in A$. Por tanto, se verifica que $\partial A \subseteq A$.

\item $\displaystyle \boxed{\supseteq}\ $ Esta inclusión es trivial, pues sabemos que $A \subseteq \overline{A}$, y por tanto $A \subseteq \partial A \cup \mathring{A} = \overline{A}$.
\end{description}

De esta forma, queda probada la equivalencia.
\end{proof}

Podemos dar ahora un teorema importante que caracteriza a los compactos de $\R^n$.

\begin{nth}[Teorema de Heine-Borel]
Sea $A\subseteq \R^n$. Entonces:
\[
A \ es \ compacto \iff \text{ A es cerrado y acotado}
\]
\end{nth}
\begin{proof}\hfill\\
	\boxed{\Rightarrow}\\
	Suponemos que $A \subset \mathbb{R}^n$ es compacto. Entonces, por su definición, $\forall \{x_n\} \subset A \ \exists \{x_{\sigma(n)}\}$ parcial de $\{x_n\}$ con $\{x_{\sigma(n)}\} \to x \in A$. Supongamos que $A$ no está acotado. Entonces, $\forall n \in \mathbb{N}, \ \exists a_n \in A : \ |a_n| \geq n$, por lo que $\{a_n\}$ no converge y por tanto $\sigma(n) \geq n \implies \{a_{\sigma(n)}\}$ no converge, por lo que $A$ está acotado.
	
	Supongamos ahora que $\{x_n\} \to x \implies \exists \{x_{\sigma(n)}\} \to x \in A$, y como sabemos que si una sucesión es convergente todas sus parciales convergen al mismo límite, entonces eso implica que $\{x_n\}\to x \in A$ por lo que toda sucesión converge a un punto de $A$, y así $A$ es cerrado.
	\\
	\boxed{\Leftarrow}\\
	Supongamos ahora que $A$ es cerrado y acotado. Sea $\{x_n\}$ una sucesión cualquiera de puntos de $A$.
	
	Como $A$ es acotado, entonces $\exists R > 0 : \ \ A \subset B(0,R)$. Además, como $\{x_n\} \subset A \ \forall n \implies |x_n| < R \ \forall n \in \mathbb{N}$, así $\{x_n\}$ es acotada.
	
	Como $\{x_n\}$ es acotada, por el teorema de Bolzano Weierstrass, $\exists \sigma: \mathbb{N} \to \mathbb{N}$ estrictamente creciente con $\{x_{\sigma(n)}\}\to x \in \mathbb{R}^n$, y como $\{x_{\sigma(n)}\}$ es una subsucesión de puntos de $A$ que converge a $x$ y el conjunto $A$ es cerrado, entonces el límite de esta sucesión está en $A$ , es decir:
	\[
	\{x_{\sigma(n)}\}\to x \in A
	\]
Por lo que tenemos la definición de conjunto compacto.
\end{proof}

\begin{nprop}
Sea $\{x_n\} \subseteq \mathbb{R}^N$ convergente a un $x_o \in \mathbb{R}^N$. Entonces, el conjunto \hfill \\$A = \{x_n:\ n=0,1,2,\dots \}$ es compacto.
\end{nprop}
\begin{proof}
	La prueba de este resultado es sencilla. Como $\{x_n\}$ es convergente, entonces está acotada. Además, el conjunto de los puntos de la sucesión es un conjunto cerrado pues los puntos son conjuntos cerrados. Así, el conjunto $A$ es cerrado y acotado, luego es compacto.
\end{proof}



%%% Clasificación de conjuntos en R^N.

\subsection{Clasificación de conjuntos en $\mathbb{R}^N$}
En esta sección, vamos a intentar definir una serie de tipos de conjuntos sobre los que luego trataremos de establecer ciertas propiedades
\begin{ndef}[Conjunto convexo]
Un conjunto $A\subseteq \mathbb{R}^N$ se dice \textit{convexo} si $\forall x,y \in A$ se tiene que el segmento de extremos $x$ e $y$ está incluido en $A$. En otras palabras: $$A\ convexo\ \iff [x,y] = \{tx + (1-t)y: \ t\in [0,1]\} \subseteq A.$$
\end{ndef}




\begin{ndef}[Poligonalmente convexo]
Un conjunto $A\subseteq \mathbb{R}^N$ se dice \textit{poligonalmente convexo} si  $\forall x,y \in A$ existe una poligonal que los une y no se sale de $A$. En otras palabras:  $A\ poligonalmente\ convexo \iff \exists \{x= a_0, a_1,\dots,a_k=y \}\subseteq A$ tal que: $$\bigcup_{i=1}^k [a_{i-1},a_i] \subseteq A.$$
\end{ndef}



\begin{ndef}[Conjunto arco-conexo]
Un conjunto $A \subseteq \mathbb{R}^N$ se dice \emph{arco-conexo(conexo por arcos)} si $\forall x,y \in A$ existe un camino incluido en $A$ que los une. En otras palabras,  $A\ es\ conexo\ por\ arcos \iff \exists \varphi:[a,b] \longrightarrow \mathbb{R}^N$ verificando: $$\varphi(a) = x;\quad \varphi(b) = y;\quad \varphi([a,b]) \subseteq A.$$
\end{ndef}



\begin{ndef}[Conjunto no conexo]
Decimos que un conjunto $A\in \mathbb{R}^N$ es \textit{NO conexo} si existen $U,\ V$ abiertos en $\mathbb{R}^N$ tales que: $$U \cap A \ne \emptyset;\quad V \cap A \ne \emptyset;\quad A \subseteq U \cup V;\quad A \cap U \cap V = \emptyset.$$
\end{ndef}



\begin{nota}
La misma definición se aplica para un espacio topológico $(X,\tau).$
\end{nota}



\begin{ndef}[Conjunto conexo]
Un conjunto $A\subseteq \mathbb{R}^N$ se dice conexo si no es no conexo. Equivalentemente, $\forall \ U,V$ abiertos en $\mathbb{R}^N$ tales que $U \cap A \ne \emptyset, \ V \cap A \ne \emptyset,\ A \subseteq U \cup V,$ se tiene que forzosamente $A \cap U \cap V \ne \emptyset$.
\end{ndef}



\begin{nprop}
Sea $A\subseteq \mathbb{R}$ un conjunto arco-conexo. Entonces, $A$ es convexo.
\end{nprop}

\begin{proof}
Sean $x,y \in A$, y supongamos sin pérdida de generalidad que $x \le y$. Sabemos que por ser $A$ arco-conexo, $\exists \varphi : [a,b] \longrightarrow \mathbb{R}$ función continua verificando: $$\varphi(a) = x;\quad \varphi(b) = y;\quad \varphi([a,b]) \subseteq A.$$
Como $\varphi$ es una función continua definida en un intervalo cerrado y acotado, aplicamos el \textbf{teorema del valor intermedio} en $\mathbb{R}$, y obtenemos que $\varphi([a,b])$ es un intervalo. Por ser un intervalo, verificará que $\forall \alpha, \beta \in \varphi([a,b])\ con\ \alpha \le \beta,\ \ [\alpha,\beta] \subseteq \varphi([a,b])$.

Por tanto, como $\varphi(a), \varphi(b) \in \varphi([a,b])$, concluimos que: $$[\varphi(a),\ \varphi(b)] = [x,y] \subseteq \varphi([a,b]) \subseteq A.$$
Así, hemos demostrado que $\forall x,y \in A\ \ [x,y] \subseteq A$, y por tanto, $A$ es convexo.
\end{proof}



\begin{nprop}
Sea $A \subseteq \mathbb{R}^N$ convexo. Entonces, $A$ es arco-conexo.
\end{nprop}

\begin{proof} \hfill \\
Fijemos $x,y \in A$ arbitrarios, y construyamos la aplicación $\varphi: [0,1] \longrightarrow \mathbb{R}^N$ dada por: $$\varphi(t) = (1-t)x + ty\ \ \forall t\in [0,1]$$
Una primera observación es que $\varphi([0,1]) = \{(1-t)x + ty: t \in [0,1]\} = [x,y] \subseteq A$ por ser $A$ convexo. También se desprende de la definición de $\varphi$ que $\varphi(0) = x$ y $\varphi(1) = y$.

Para comprobar que $\varphi$ es continua, utilicemos la caracterización de la continuidad por sucesiones:

Sea $\{x_n\} \subseteq [0,1]\ con\ \{x_n\} \rightarrow a \in [0,1]$. Entonces, $\{\varphi(x_n) \} = \{(1- x_n)x + x_n y\}$. \\Apliquemos ahora propiedades de las sucesiones convergentes, y obtenemos que: $$\{\varphi(x_n) \} \rightarrow (1-a)x + ay = \varphi(a).$$
Entonces, $\forall \{x_n\} \subseteq [0,1]\ con\ \{x_n\} \rightarrow a \Rightarrow \{\varphi(x_n) \} \rightarrow \varphi(a)$, por lo que $\varphi$ es continua.

Así, queda probado que $A$ es conexo por arcos.
\end{proof}



\begin{nprop}
Sea $A\in \mathbb{R}^N$ arco-conexo. Entonces, $A$ es conexo.
\end{nprop}



\begin{nprop}
Sea $A\subseteq \mathbb{R}^N$ abierto y conexo por arcos. Entonces, $A$ es poligonalmente convexo.
\end{nprop}


%%% Continuidad en espacios topológicos. Topología inducida.
\subsection{Continuidad en espacios topológicos. Topología inducida.}

\begin{ndef}[Continuidad en espacios topológicos]
Sean $(X,\tau_x),\ (Y,\tau_y)$ dos espacios topológicos, y sea $f:X\longrightarrow Y$. Entonces: $$f\ es\ continua \iff f^{-1}(B) \in \tau_x \ \ \forall B \in \tau_y.$$
\end{ndef}



\begin{ndef}[Topología inducida]
Sea $(X,\tau)$ un espacio topológico, y $A\subseteq X$. Entonces, $\tau_A = \{B\cap A: \ B \in \tau\}$ es la \textit{topología inducida en $A$}.
\end{ndef}



\begin{nprop}[Caracterización de abiertos en topología inducida]
Sea $(X,\tau)$ un espacio topológico, y $A\subseteq X$. Si $(A,\tau_A)$ es el espacio topológico inducido en $A$, entonces: $$B' \in \tau_A \iff \exists B\in \tau: \ B' = B\cap A.$$
\end{nprop}



\begin{nprop}
Sea $(X,\tau)$ un espacio topológico, y $A\subseteq X$. Entonces, $A$ es no conexo si, y solo si, existen $U,V$ \textbf{abiertos en $\bm{(A,\tau_A)}$} tales que: $$U \ne \emptyset \ne V;\quad A \subseteq U \cup V;\quad U \cap V = \emptyset.$$
\end{nprop}



\begin{ndef}[Continuidad en topología inducida]
Sean $(X,\tau_x),\ (Y,\tau_y)$ dos espacios topológicos, $A\subseteq X$, y $f:A\longrightarrow Y$. Entonces: $$f\ es\ continua \iff f\ es\ continua\ en\ (A,\tau_A).$$
\end{ndef}



%%% Teoremas sobre funciones continuas en R^N.
\subsection{Teoremas sobre funciones continuas en $\mathbb{R}^N$}

A continuación, daremos ciertos de los teoremas más usados en el análisis elemental. Se revisitarán los teoremas dados con anterioridad para funciones continuas en $\R$ y se enunciarán y probarán sus versiones correspondientes en $\R^n$ o en espacios métricos en general.



\begin{nth}[Weierstrass]
Sea $(X,d)$ un espacio métrico, $\emptyset \ne A \subseteq X$ compacto, y $f:A \longrightarrow \mathbb{R}$ continua en $A$. Entonces, $\exists x_1,x_2 \in A: \ f(x_1)\le f(x) \le f(x_2)\ \ \forall x\in A$. En otras palabras, la función $f$ alcanza su mínimo y su máximo.
\end{nth}

\begin{nth}[Weierstrass generalizado]
Sean $(X,d)$, $(Y,d)$ espacios métricos, $\emptyset \ne K \subseteq X$ compacto, y $f: K \longrightarrow Y$ continua. Entonces, $f(K)$ es compacto.
\end{nth}
\begin{proof}
		Sea $\{y_n\}\subseteq f(K)$ una sucesión cualquiera en $f(K)$. Para demostrar que $f(K)$ es compacto debemos demostrar que $\{y_n\}$ tiene una subsucesión convergente a algún punto de $f(K)$. Sea $y_n$ = $f(x_n)$ con $x_n \in K$, por ser A compacto $\{x_n\}$ tiene una subsucesión $\{x_{\sigma(n)}\}$ convergente a un $a \in K$. Por lo tanto, $\{f(x_{\sigma(n)})\}$ convergente a un $f(a) \in f(K)$, siendo $\{f(x_{\sigma(n)})\}$ una subsucesión de $\{f(x_{n})\}$, es decir, de $\{y_n\}$.
\end{proof}



\begin{nth}[Valor Intermedio]
Sea $\emptyset \ne A \subseteq \mathbb{R}^N$ arco conexo, y $f: A \longrightarrow \mathbb{R}^M$ continua. Entonces, $f(A)$ es arco-conexo en $\mathbb{R}^M$.
\end{nth}

\begin{proof}
Sean $X,Y\in f(A)$. Entonces, $\exists x,y \in A : X=f(x), \ Y=f(y)$. Como $A$ es arco-conexo, $\exists\varphi : [a,b]\longrightarrow \mathbb{R}^N \text{ continua tal que}\ \varphi(a) = x,\; \varphi(b)=y,\; \varphi([a,b]) \subseteq A$.


Ahora, definimos $\psi := f\circ \varphi : [a,b] \longrightarrow \mathbb{R}^M$, que es continua por ser composición de funciones continuas. Entonces, se verifica que: $$\psi(a) = f(\varphi(a)) = f(x) = X;\quad \psi(b)= f(\varphi(b)) = f(y) = Y;\quad \psi([a,b]) = f(\varphi([a,b])) \subseteq f(A).$$
Por tanto, queda probado que $f(A)$ es arco-conexo en $\mathbb{R}^M$.
\end{proof}



\begin{nth}[Valor Intermedio revisitado]
Sea $\emptyset \ne A\subseteq \mathbb{R}^N$ conexo, y $f:A \longrightarrow \mathbb{R}^M$ continua. Entonces, $f(A)$ es conexo en $\mathbb{R}^M$.
\end{nth}

El siguiente resultado, tiene mucha relevancia en el estudio de la continuidad uniforme de funciones en un espacio métrico como es $\R^n$.
\begin{nth}[Heine-Cantor]
	Sea $\emptyset\ne A \subseteq \mathbb{R}^N$ compacto, y $f : A \longrightarrow \mathbb{R}^N$ continua. Entonces $f$ es uniformemente continua en $A$.
\end{nth}


\begin{proof}
	$f$ es continua en A $\implies$ $f$ es continua en $a\;\; \forall a \in A$. Ahora, sea $\epsilon>0$ fijo.
	\[
	    \forall a \in A \quad \exists \delta = \delta_a > 0\;\; \forall x\in A\;\; d(x,a) < \delta_a \implies d(f(x),f(a))<\epsilon
	\]

	Tomamos un recubrimiento abierto de $A$, y como $A$ es compacto, encontramos un subrecubrimiento finito.
	
	\[
		A \subseteq \bigcup_{a\in A} B(a, \frac{\delta_a}{2}) \implies \exists a_1,\dots,a_n \in A: A \subseteq \bigcup_{i=1}^n B\left(a_i, \frac{\delta_{a_i}}{2}\right)
	\]
	
	Por esta última inclusión:
	\[
		\forall x\in A\quad \exists i \in \left \{ 1,\dots,n \right \} : x\in B\left(a_i,\frac{\delta_{a_i}}{2}\right)\cap A \implies f(x)\in B(f(a_i),\epsilon)
	\]
	
	Sean $\delta = \min\left\{\ddfrac{\delta_{a_i}}{2} : i \in \left \{ 1,\dots,n \right \}\right\} > 0$ y $y\in A : d(x,y) < \delta < \delta_{a_i}$ para un $x\in A$ fijo. Tomamos el $a_i$ proporcionado por la proposición anterior para $x$.
	\[
		d(y,a_i) \le d(y,x)+d(x,a_i) < \delta_{a_i} \implies y\in B(a_i,\delta_{a_i}) \implies f(y) \in B(f(a_i), \epsilon)
	\]
	
	Finalmente,
	\[
		d(f(x), f(y)) \le d(f(x),f(a_i)) + d(f(a_i), f(y)) < \epsilon
	\]
	
	Para cualquier $\epsilon$ para el que se desee que se verifique la condición de la continuidad uniforme, basta tomar $\ddfrac{\epsilon}{2}$ en la continuidad.
\end{proof}

\begin{proof}[Demostración alternativa]
	La condición para la continuidad uniforme es la siguiente:
	\[
		\forall \varepsilon > 0 \  \exists \delta > 0 \ \forall x, y \in A : d(x,y) < \delta \implies d (f(x) , f(y) ) < \varepsilon
	\]
	
	Vamos a proceder por reducción al absurdo, para lo cual negamos esta condición:
	\[
		\exists \varepsilon_0 > 0 \ \forall \delta > 0 \ \exists x, y \in A : d (x,y) < \delta \wedge d (f(x) , f(y) ) \ge \varepsilon_0
	\]
	
	Tomamos este $\epsilon_0$, lo que nos da, para cada $\delta>0$, un par de puntos $x$ e $y$ que cumplen la propiedad expresada arriba. Tomamos $\delta = \frac{1}{n} \ \forall n\in \mathbb{N}$. Esto nos da dos sucesiones $\{x_n\}$ e $\{y_n\}$ tales que
	\[
		d(x_n,y_n) < \frac{1}{n} \wedge d(f(x_n),f(y_n)) \ge \epsilon_0
	\]
	
	Por ser A compacto, el teorema de Bolzano-Weierstrass nos da dos sucesiones parciales $\{x_{n_k}\}$ a $x_0$ e $\{y_{n_k}\}$ a $y_0$. Por tanto:
	\[
		d(x_{n_k},y_{n_k}) < \frac{1}{n_k} \wedge d(f(x_{n_k}),f(y_{n_k})) \ge \epsilon_0
	\]
	
	Sin embargo, $\{x_{n_k}\}$ e $\{y_{n_k}\}$ convergen al mismo punto (por converger su distancia a cero), y como $f$ es continua, esta proposición no puede ser verdadera. Hemos llegado por tanto a una contradicción, luego $f$ debe ser uniformemente continua.
\end{proof}


\section{Límite funcional en $\mathbb{R}^N$.}

\begin{ndef}[Límite funcional]
	Sean $\emptyset \ne A \subseteq \mathbb{R}^N$, $a\in A'$ y $f: A \longrightarrow \mathbb{R}^M$. Entonces $f$ tiene límite $l$ en $x=a$, y se denota $\displaystyle\lim_{x\to a} f(x)$ si:
	
\[
	\forall \epsilon>0\quad \exists \delta>0 : \begin{rcases}
	0<d(x,a)<\delta\\
	x\in A
\end{rcases} \implies d(f(x), l) < \epsilon
\]
\end{ndef}



\begin{nprop}[Caracterización punto de acumulación]
Sea $(X,d)$ un espacio métrico, y $A\subseteq X$. Consideremos un punto $x\in X$. Son equivalentes:

\begin{nlist}
\item x es punto de acumulación de A.
\item $\exists \{a_n\}\subseteq A-\{x\}$ tal que $\{a_n\} \rightarrow x$.
\item $\forall \epsilon > 0\ B(x,\epsilon)\cap (A-\{x\})$ es un conjunto infinito. 
\end{nlist}
\end{nprop}



\section{Funciones derivables en $\mathbb{R}^N$.}
\subsection{Concepto de función derivable.}

Sea $A$ un abierto de $\mathbb{R}^N$. Partimos de la siguiente observación:
\[
	\forall x_0 \in A\quad \exists \delta >0 : B(x_0, \delta) \subset A \implies \forall v \in \mathbb{R}^N\quad \exists \epsilon > 0 : [\; t \in (-\epsilon, \epsilon) \implies x_0+tv\in B(x_0,\delta)\;]
\]
 En particular, si $|v| = 1 \implies \epsilon = \delta$.
 
\begin{ndef}[Función derivable] Sean $f : A \longrightarrow \mathbb{R}^M$ y $x_0\in A$. Se dice que $f$ es derivable en $x_0$, según Fréchet, si
\[
	\exists L\in Lin(\mathbb{R}^N, \mathbb{R}^M) : \lim_{x\to x_0} \frac{|f(x)-f(x_0)- L(x-x_0)|}{|x-x_0|} =0
\]

Notamos $Df(x_0) = L$.
\end{ndef}

\begin{nota}[1]\hfill
	\begin{nlist}
		\item El límite tiene sentido porque $x_0\in A$.
		\item El límite anterior es equivalente a $\lim_{y\to 0} \frac{|f(x_0+y)-f(x_0)-L(y)|}{|y|}$.
	\end{nlist}
\end{nota}

\begin{nota}[2]
	$A$ es abierto $\implies$ $L$ (si existe) es única. De aquí se exige que $A$ sea un abierto. 
\end{nota}

\begin{proof}[Demostración (Nota 2)]
	Suponemos que $\exists L_1,L_2\in Lin(\mathbb{R}^N, \mathbb{R}^M)$ tales que 
\[
\lim_{x\to x_0} \frac{|f(x)-f(x_0)- L_1(x-x_0)|}{|x-x_0|} =0= \frac{|f(x)-f(x_0)- L_2(x-x_0)|}{|x-x_0|}
\]

Entonces, dado un $x\in A$:

\begin{align}
	\frac{|L_1(x-x_0) - L_2(x-x_0)|}{|x-x_0|} \le \frac{|f(x)-f(x_0)- L_1(x-x_0)|}{|x-x_0|} + \frac{|f(x)-f(x_0)- L_2(x-x_0)|}{|x-x_0|}\notag
\end{align}

\[
\implies \lim_{x\to x_0} \frac{|(L_1-L_2)(x-x_0)|}{|x-x_0|} = 0 \implies \lim_{x\to x_0}|(L_1 -L_2)(x-x_0)|=0
\]
Como A es abierto, sea $y \in \mathbb{R}^N, y \neq 0 \implies x:=x_0+y \in B(x_0, \delta) \subset A$\\
$$\lim_{x\to x_0}(L_1-L_2)(x-x_0)=0=\lim_{y\to 0}(L_1-L_2)(y)\implies L_1=L_2$$
\end{proof}

\begin{nprop} En los mismos términos: si $f$ es derivable en $x_0$ $\implies$ $f$ es continua en $x_0$.
	
\end{nprop}

\begin{proof} Para probar esta proposición, hay que probar que $\lim_{x\to x_0} (f(x)-f(x_0)) = 0$
\[
	\lim_{x\to x_0} (f(x)-f(x_0)) = \underbrace{\lim_{x\to x_0} (f(x)-f(x_0)-L(x-x_0))}_{=\;0 \text{ ($f$ derivable)}} + \underbrace{\lim_{x-x_0} L(x-x_0)}_{=\; 0 \text{ ($L$ lineal $\implies$ continua)}} = 0
\]
\end{proof}

\begin{ndef}[Derivada direccional]
	Sea $v\in \mathbb{R}^N$, con $|v| = 1$. $f$ es derivable en $x_0$ en la dirección $v$ si:
	\[
		\exists \lim_{t\to 0} \frac{f(x_0+tv)-f(x_0)}{t} = D_v f(x_0) \iff
	\]
	
	\begin{center}
	$\iff f_1,f_2,\cdots f_m$ derivable direccionalmente en $x_0$ en la dirección v. $\iff$
	\[
	\iff D_vf(x_0) = (D_vf_1(x_0),\cdots,D_vf_m(x_0))
	\]
\end{center}
\end{ndef}


\begin{nprop}
	Sea $f$ derivable en $x_0\implies f$ derivable a lo largo de la dirección v y $D_vf(x_0) = Df(x_0)(v)$
	\begin{proof}
	$f$ derivable en $x_0$. Tomo $y=tv  $. Podemos ver entonces:
	
	\[
	\lim_{t \to 0} \frac{f(x_0 +tv) - f(x_0) - Df(x_0)(tv)}{t} = 0 \implies
	\]
	\[
	\implies \lim_{t \to 0} \abs{\frac{f(x_0 +tv) - f(x_0)}{t} - Df(x_0)(\frac{tv}{t})} \implies \lim_{t \to 0}\frac{f(x_0 +tv) -f(x_0)}{t} - Df(x_0)(v)
	\]
	\[
	\implies \lim_{t \to 0}\frac{f(x_0 +tv) - f(x_0)}{t} = Df(x_0)(v)
	\]
Hemos probado que $\exists D_vf(x_0)$
\end{proof}
\end{nprop}


%% Introducir definición alternativa (usando la definición de límite directamente).
\section{Matriz asociada a $Df(x_0)$.}
	Si $f : B \longrightarrow \mathbb{R}^M$, con $\emptyset \ne B \in \mathbb{R}^N$, la matriz asociada a $Df(x_0)$ es una matriz de orden $M \times N$, (que notaremos $A$ en lo sucesivo), como ya sabemos, por ser una aplicación lineal. Ahora, nuestro siguiente objetivo es encontrar esa matriz. Observemos cómo podemos obtenerla por filas, aplicándole los vectores de la base canónica:
	
	\[
		e_i = (0, \dots, \stackrel{i)}{1}, \dots, 0) \implies Df(x_0)(e_i) = D_{e_i}f(x_0) = Ae_i = \begin{pmatrix}
	a_{1i}\\ \vdots\\ a_{Mi}
\end{pmatrix}
	\]

Tras esta observación, vamos a caracterizar cada elemento $a_{ji}$:

\begin{align*}
	D_{e_i}f(x_0) = \\\lim_{t\to 0} \frac{f(x_0+te_i)-f(x_0)}{t} = \\\lim_{t\to 0} \left(\frac{f_1(x_0+(0,\dots,\stackrel{i)}{t},\dots,0)-f_1(x_0)}{t},\right.\left.\cdots, \frac{f_M(x_0+(0,\dots,\stackrel{i)}{t},\dots,0)-f_M(x_0)}{t}\right)=\\
	 \left(\lim_{t\to 0} \frac{f_1(x_0+(0,\dots,\stackrel{i)}{t},\dots,0)-f_1(x_0)}{t}, \cdots, \lim_{t\to 0} \frac{f_M(x_0+(0,\dots,\stackrel{i)}{t},\dots,0)-f_M(x_0)}{t}\right) =\\
	 \left(D_{e_i}f_1(x_0, \dots, D_{e_i}f_M(x_0))\right) = \left(\frac{\partial f_1(x_0)}{\partial x_i}, \dots, \frac{\partial f_M(x_0)}{\partial x_i}\right) \implies a_{ji} = \frac{\partial f_j}{\partial x_i} (x_0)
\end{align*}

Por tanto, $A$ queda de la siguiente forma:

\[
	A = \begin{pmatrix}
	\frac{\partial f_1(x_0)}{\partial x_1} (x_0) & \cdots & \frac{\partial f_1(x_0)}{\partial x_N} (x_0)\\
	\vdots & \ddots & \vdots\\
	\frac{\partial f_M(x_0)}{\partial x_1} (x_0) & \cdots & \frac{\partial f_M(x_0)}{\partial x_N} (x_0)
\end{pmatrix}
\]

Deducimos que:

\[
	\exists \lim_{t\to 0} \frac{f(x_0+tv) - f(x_0)}{t} := D_vf(x_0) \iff f_1, \dots, f_M \text{ son derivables direc. en $x_0$ en la dir. $v$.}
\]

Además, $D_vf(x_0) = (D_vf_1(x_0), \dots, D_vf_M(x_0))$.







\section{FALTA MUCHO CONTENIDO}

\begin{nprop}
	Sea $L$ una aplicación lineal. Entonces L es derivable y $DL(a) = L$
	
\end{nprop}
\begin{proof}
	
	Tomaremos:
	\[
	lim_{x\to a} \frac{||Lx -La - M(x-a)||}{||x-a||}= 0
	\]
	Entonces, tenemos que encontrar : $M: R^n \to R^m$ lineal para que sea derivable. Entonces, podemos tomar $L=M$ y como $L$ es lineal, es trivial que es verdad (por la linealidad de L) y esto implica que $L$ es derivable y $DL(a) = M = L$. 
\end{proof}

\begin{nprop}
	Toda aplicación lineal es Lipschitziana. Si es Lipschitziana entonces es continua.
\end{nprop}
\begin{proof}
	Como $L$ es lineal, entonces L es derivable y por tanto continua. Entonces:
	\[
	||Lx|| \leq M||x|| \quad \forall x \in R^n	
	\]
	\[
	 ||Lx-Ly|| = ||L(x-y) || \leq M||x-y|| \implies L\ es \ Lipschitziana
	\]
\end{proof}


\begin{nprop}
	Si $L:R^n \to R^m$ lineal $\implies L $ es continua y $\exists M \geq 0: ||Lx||_{R^m} \leq M ||x||_{R^n} \quad \forall x \in R^n$
\end{nprop}
\begin{proof}
	Si tomamos $L|_{S(0,1)}$ y vemos que es acotada, entonces podríamos definir: \[M = sup\{Lx : x \in S(0,1)\} \implies ||Lx|| \leq M \quad \forall ||x|| = 1 \implies ||L(x)|| = ||\ ||x||\ L(\frac{x}{||x||}) ||\] donde hemos usado que $L$ es lineal. Y esto implica: \[||\ ||x||\ || L(\frac{x}{||x||})|| || = ||x||\ || L(\frac{x}{||x||}) || \leq M||x|| \quad \forall x \in R^n-\{0\}\]
\end{proof}


\begin{nprop}
	Si $\begin{rcases}
	B: R^N x R^n \to R\\
	Bilineal
\end{rcases} B \ es \ continua$.

Es más, $\exists M \geq 0 : |B(x,y)| \leq M||x||\||y|| \quad \forall x,y \in R^n$
	
	Se deja como ejercicio la demostración. Se debe usar que :
	
	\[
	B(x,y) = ||x||\ ||y|| B(\frac{x}{||x||},\frac{y}{||y||}) 
	\]
	
	Y tomar el $M = sup\{|B(x,y): ||x|| = 1 \ y \ ||y|| = 1\}$ y para ello necesito ver que $S_{R^n}(0,1)xS_{R^n}(0,1)\{(x,y) \in R^nxR^n : ||x|| = 1 \ y \ ||y|| = 1\}$ es compacto.
	
	Quién sería en este caso el candidato a $DB(x_0,y_0)(u,v) = D(x_0,v)+B(u,y_0)$.
	
\[
\lim_{(x,y) \to (x_0,y_0)}\frac{|B(x,y) - B(x_0,y_0) - B(x_0,y-y_0) + B(x-x_0,y_0)}{||(x-x_0,y-y_0)||}
\]
\[
\lim_{(x,y) \to (x_0,y_0)}\frac{|B(x,y) - B(x_0,y_0) - B(x_0,y) + B(x_0,y_0) - B(x,y_0) + B(x_0,y_0)|}{||(x-x_0,y-y_0)||}
\]
\[
\lim_{(x,y) \to (x_0,y_0)}\frac{|B(x-x_0,y)- B(x-x_0,y_0)|}{||(x-x_0,y-y_0)||} = \lim_{(x,y) \to (x_0,y_0)}\frac{|B(x-x_0,y-y_0)|}{||(x-x_0,y-y_0)||}
\]

Y ahora, como $(x,y) \to (x_0,y_0) \implies (x-x_0,y-y_0) \to (0,0)$ y ahora si ponemos $(x-x_0,y-y_0) = (u,v)$

\[
 \lim_{(u,v) \to (0,0)}\frac{|B(u,v)|}{||(u,v)||} = 0
\]

Estudiando ese límite por cualquier método, nos saldría 0. Luego, tenemos que intentarlo pasando el problema a coordenadas polares o con algún otro método.
\end{nprop}

\begin{ejemplo}
	Sea $f:R^2 \to R$ bilineal con $f(x,y) = xy$.
	
	Entonces, en un punto $(x_0,y_0)$:
	
	\[
	\lim_{(x,y) \to (x_0, y_0)} \frac{|x(x,y) - f(x_0,y_0) - L(x-x_0, y-y_0)|}{||(x-x_0,y-y_0||} = 
	\]
	\[
	 =  \lim_{(x,y) \to (x_0,y_0)} \frac{|xy-x_0y_0 - L(x-x_0,y-y_0)}{||(x-x_0,y-y_0||}
	\]
	
	Pero, ¿cuál es esa $L$?. Si tomáramos $L(x-x_0,y-y_0) = y_0(x-x0) - x_0(y-y_0)$ en el numerador nos quedaría $|xy - x_0y_0 - y_0(x-x0) + x_0(y-y_0)| = |xy- x_0y_0 -y_0x -x_0y_0 - x_0y+x_0y_0 |= | (x-x_0)(y-y_0)|$. Y tenemos en el límite que:
	\[
	 = \lim_{(x,y) \to (x_0,y_0)} \frac{|(x-x_0)(y-y_0)|}{||(x-x_0,y-y_0))||} = 0
	\]
	Así, la derivada $Df(x_0,y_0)(u,v) = x_0v + y_0 u$
\end{ejemplo}

\begin{nprop}
	Sea $g,f: R^n \to R^m$ derivable en $a \in R^n$. El producto escalar es derivable.
	
\begin{proof}
		Construyo $h: R^n \to R$ con $h(x) = < f(x), g(x)>$. Probar que es derivable:
	
	$h = B o (f,g)$ con $B(x,y) = <x,y> \quad \forall x,y \in R^n$
	
	La derivada de la h es:
	\[
	Dh(x_0)(x) = [DB(f(x_0)g(x_0) o D(f,g)(x_0)](x) =\]\[DB(f(x_0),g(x_0))[D(f,g)(x_0)(x)] = DB(f(x_0),g(x_0))*[Jacob(f)](x) = 
	\]
	\[
	= B(f(x_0),[Jacob(g)(x)]) + B([Jacob(f)(x), g(x_0)] = \]
	\[
	 = B(f(x_0),Dg(x_0)(x)) + B(Df(x_0)(x),g(x_0)) = 
	\]
	\[
	<f(x_0),Dg(x_0)(x) > + < Df(x_0)(x),g(x_0)>
	\]
\end{proof}
	
\end{nprop}

\begin{nota}
	La aplicación $R^n \to R$ que se lleva $y \to y^t \ \mathcal{H}f(c)y$  Forma cuadrática asociada a la aplicación bilineal $B:R^nxR^n \to R$ tal que $(y,z) \to y^t\mathcal{H}f(c)z$ la matriz Hessiana $\mathcal{H}f(c)$ es simétrica. Se puede pruede probar que:


\end{nota}

\textbf{Lema de Schwarz}
Sea $A\subset R^n$ no vacío, y $f$ una función de clase $\varphi^2(A)$
$\forall i,j = 1,...,n$ con $i\ne j$
\[
\frac{\partial^2f}{\partial x_i \partial x_j}(x) = \frac{\partial^2 f}{\partial x_i \partial x_i}(x)
\]
\begin{proof}

Vamos a hacer esta prueba para $N=2$.

Si $(x,y) \in A \implies \exists h_0 \ , \ k_0 > 0 : \begin{rcases}
	0 < h < h_0 \\
	0 < k < k_0
\end{rcases} \implies (x,y),(x+h,y),(x,y+k),(x+h,y+k) \in A$

Podemos escribir:
\[
S_{k,h}(x) =[f(x+h,y+k) - f(x,y+k)] - [f(x+h,y) -f(x,y)]
\]
	
Como $f$ es de clase 2, es de clase 1 y podemos usar el TVM para evaluar cuánto valen las diferencias que hemos expresado en $S_{k,h}(x)$. Esa diferencia se puede estimar como la derivada en el segmento aplicado al vector diferencia de esos dos puntos.
\[
TVM \implies Df(c)[(x+h-x,(y+k)-(y-k)] -Df(d)[(x+h-x,y-y)]= \frac{\partial f}{\partial x}(c)h - \frac{\partial f}{\partial x}(d)h
\]
Y podemos volver a usar el TVM pues sigue siendo de clase 1.

Por otro lado, vemos que:
\[
S_{k,h}(x)= [f(x+h,y+k)-f(x+h,y)]-[f(x,y+k)-f(x,y)]
\]
Usando de nuevo el TVM obtenemos:
\[
\frac{\partial f}{\partial y}(e)k - \frac{\partial f}{\partial y}(g)k
\]

El teorema sin embargo no nos dice dónde están esos puntos y (por estar en $R^2$) tenemos que buscar que los puntos que nos de el TVM estén alineados.

Supongamos que $c , d$ están siempre en la misma vertical. Entonces, usando el TVM para $\frac{\partial f}{\partial x}$:
\[
\frac{\partial}{\partial y}(\frac{\partial f}{\partial x})(P)kh
\]

Y supongamos que $e,g$ están en la misma vertical. Usando el TVM para $\frac{\partial f}{\partial y}$
\[
\frac{\partial }{\partial x}(\frac{\partial f}{\partial y})(Q)hk
\]
Como ambos eran iguales, entonces:
\[
\frac{\partial}{\partial y}(\frac{\partial f}{\partial x})(P)kh = \frac{\partial }{\partial x}(\frac{\partial f}{\partial y})(Q)hk \implies \frac{\partial}{\partial y}(\frac{\partial f}{\partial x})(P)=\frac{\partial }{\partial x}(\frac{\partial f}{\partial y})(Q)
\]
Si lo que hacemos es que $h \to 0 \implies x+h\to x$ y lo mismo con la $k$ y así $P,Q\to (x,y)$ y de esta forma:
\[
\frac{\partial}{\partial y}(\frac{\partial f}{\partial x})(P) \to \frac{\partial^2 f}{\partial x \partial y}(x,y)
\]
\[
\frac{\partial}{\partial x}(\frac{\partial f}{\partial y})(Q) \to \frac{\partial^2 f}{\partial y \partial x}(x,y)
\]
Y así, como eran iguales y tienden a lo mismo, tenemos que:
\[
\frac{\partial^2 f}{\partial x \partial y}(x,y) =  \frac{\partial^2 f}{\partial y \partial x}(x,y)
\]

Pero habíamos supuesto que c y d están siempre en la misma vertical. Ahora, lo que tendríamos que hacer es intentar haber usado sólamente una vez el TVM.
Par ello, definimos:
\[
F_h(k) = f(x+h,y+k)-f(x,y+k)
\]

Con esta definición, tenemos $F_h(0) = f(x+h,y)-f(x,y)$ y ahora si hacemos:
\[
S_{k,h}(x) =[f(x+h,y+k) - f(x,y+k)] - [f(x+h,y) -f(x,y) = F_h(k) - F_h(0)
\]

Vamos a ver primero que
\[
F_h'(K) = [\frac{\partial f}{\partial x}(x+h,y+K)*0 + \frac{\partial f}{\partial y}(x+h,y+K)*1] - 
\]
\[
\frac{\partial f}{\partial x}(x,y+K)*0 + \frac{\partial f}{ \partial y}(x,y+K)*1] = \frac{\partial f}{\partial y}(x+h,y+K) - \frac{\partial f}{\partial y}(x,y+K)
\]
Luego:
\textbf{NO PUEDO COPIAR MÁS, COMPLETAR}

\end{proof}
\begin{proof}[Demostración para el caso general]
	Tenemos $f:A \to R$ y $x=(x_1,...,x_i,...,x_j,...,x_n) \mapsto f(x_1,...,x_i,...,x_j,...,x_n)$.
	
	Definimos $g(s,t) = f(x_1,...,x_{i-i},s,...,x_{j-1},t,...,x_n)$
	Si ahora hacemos
	\[
	\frac{\partial g}{\partial s \partial t}(x_i,x_j) = \frac{\partial^2 f}{\partial x_i \partial x_j}
	\]
	Y también:
		\[
	\frac{\partial g}{\partial t \partial s}(x_i,x_j) = \frac{\partial^2 f}{\partial x_j \partial x_i}
	\]
	
	Y como lo teníamos probado para $N=2$, entonces ya tenemos que las dos son iguales pues $g$ es una función de $N=2$ variables.
	

\end{proof}

\begin{nota}[Observación 1]. Si tenemos $f \in \varphi^k$, podríamos hacer:
\[
\frac{\partial^3 f}{\partial x_1 \partial x_2 \partial x_1} = \frac{\partial}{\partial x_1}(\frac{\partial}{\partial x_2}(\frac{\partial f}{\partial x_1})) = \frac{\partial}{\partial x_2 \partial x_1}(\frac{\partial f}{\partial x_1})
\]
Si desarrolláramos:
\[
\frac{\partial^3 f}{\partial x_1 \partial x_1 \partial x_2}
\]
Se podría ver que se llega a lo mismo e incluso podemos notarlo como:
\[
\frac{\partial^3 f}{\partial x_2^2 \partial x_1}
\]
\end{nota}
\begin{nota}
	Podemos escribir:
	\[
	D^{(\alpha_1,...,\alpha_n)}f = \frac{\partial^{|\alpha|}f}{\partial x_1^{\alpha_1}... \partial x_k^{\alpha_n}}
	\]
	Con $\alpha = (\alpha_1,...,\alpha_n) \in \mathcal{N}\cup \{0\}$. A esto se le llama multiíndice.
	
\end{nota}
\begin{nota}
	Si $v\in R^n$ con $|v| = 1 $ y $D_vf(a) = Df(a)(v)$. Así:
	\[
	D_{\frac{w}{|w|}}f(a) = (1/|w|)*Df(a)(w)
	\]\[
	|w|D_{\frac{w}{|w|}}f(a) = Df(a)(w)
	\]
\end{nota}

\begin{nth}[Teorema de Taylor(Orden 2)]

$A \subset R^n$ convexo  y abierto, $x_0 \in A$ y $f:A \to R\in \varphi^2(A)$

$\forall x \in A : [x,x_0] \subset A \implies f(x) = f(x_0) + \frac{1}{1!}Df(x_0)(x-x_0)+ \frac{1}{2!}(x-x_0) + \mathcal{H}f(x)(x-x_0)$ con $c \in [x,x_0]$ siendo este el segmento que une $x \ con \ x_0$

\end{nth}

\begin{nth}[Teorema de Taylor de orden n]
Sea $\emptyset \ne A \subset R^{n}$ abierto, $a\in A$ y $f \in \varphi^{n+1}(A)$

Entonces $\forall x \in A :$ el segmento que une $x$ con $a$ está contenido en A , se tiene que $\exists c $ en el segmento que une $x$ con $a$ tal que:
\[
f(x) = \sum_{|\alpha| = 0}^n \frac{D^\alpha f(a)}{\alpha !}(x-a)^\alpha + \sum_{|\alpha| = n+1}\frac{D^\alpha f(c)}{\alpha !}(x-a)^\alpha
\]

	
\end{nth}

\begin{nth}
Sea $\empty \ne A \subset R^n \ abierto \ , \ x_0 \in A$ con $f:A \to R$ de clase $C^2$
\begin{enumerate}
	\item $f$ tiene un maximo relativo en $x_0 \implies \begin{cases}
	Df(x_0) = 0 \iff Jf(x_0) = 0\\
	\mathcal{H}f(x_0) \ es \ semidefinida \ negativa
\end{cases}$

\item $\begin{rcases}
	Df(x_0) = 0\\
	\mathcal{H}f(x_0) \ es \ definida \ negativa
\end{rcases} \implies $ f tiene un máximo relativo
\end{enumerate}	
\end{nth}

\begin{proof} Probaremos 1 y 2: 
	\begin{enumerate}
	\item Como $f$ presenta un máximo relativo en $x_0\in A=abierto$ entonces $\exists r > 0 : B(x_0,r) \subseteq A $ y $f(x) \leq f(x_0) \ \forall x \in B(x_0,r)\subset A$. Por ello, ya sabemos que $Df(x_0) = 0$. Entonces, aplicamos Taylor $\forall x \in B(x_0,r)$ segmento que une $x$ con $x_0$ está contenido en $B(x_0,r) \subset A \implies $
	\[
	0 \geq f(x) - f(x_0) =  \frac{1}{2}(x-x_0)^t \mathcal{H}f(c)(x-x_0)
	\]
	Por tanto, podemos ver que 
	\[
	0 \geq (x-x_0)^t \mathcal{H}f(x_0)(x-x_0) \quad \forall x \in B(x_0,r)
	\]
	
	Pues como la matriz hessiana está llena de las derivadas segundas en $c$, cuando $c \to x_0$ entonces $\mathcal{H}f(c) \to \mathcal{H}f(x_0)$.
\[
	0 \geq f(x) - f(x_0) =  \frac{1}{2}(x-x_0)^t \mathcal{H}f(c)(x-x_0) \iff \mathcal{H}f(c) \ semidefinida \ negativa
	\]
	\[
	 \iff y^t\mathcal{H}f(c) \ y \leq 0 \quad \forall y \in B(0,r)
	\]
	Y como hemos dicho que cuando $c \to x_0$ entonces $\mathcal{H}f(c) \to \mathcal{H}f(x_0)$ entonces $0 \geq z^t \mathcal{H}f(x_0)z \quad \forall z \in R^n \implies 0 \geq <^t \mathcal{H}f(x_0)< \quad \forall z \in R^n$ y así la matriz Hessiana es semidefinida negativa.
	
	\item Como $f$ presenta un máximo relativo en $x_0\in A=abierto$ entonces $\exists r > 0 : B(x_0,r) \subseteq A $ y $f(x) \leq f(x_0) \ \forall x \in B(x_0,r)\subset A$. Por ello, ya sabemos que $Df(x_0) = 0$.Usando el teorema de Taylor, obtenemos que:
	\[
	f(x)-f(x_0) = \frac{1}{2!}(x-x_0)^t \mathcal{H}f(c)(x-x_0)
	\]
	La hipótesis que tenemos ahora es:
	\[
	z^t \mathcal{H}f(x_0)z < 0 \quad \forall z \in R^n-\{0\}
	\]
	Ahora, si el límite es menor estricto que cero implica que los números que se aproximen tienen que ser estrictamente menor que cero, por eso es definida negativa.
	
	Por tanto, tenemos que demostrar la implicación:
	\[
	\begin{rcases}
	\mathcal{H}f(x_0) \ definida \ negativa\\
	\mathcal{H}f(c) \to \mathcal{H}f(x_0) \ si (c \to x_0)
\end{rcases} \implies \mathcal{H}f(c) \ definida \ negativa
	\]

	Y teniendo eso, podemos ver que
	
	\[
	z^t \mathcal{H}f(c) < 0 \quad \forall z \in R^n-\{0\} \iff \forall x \in B(x_0,r)-\{x_0\}
	\]
	Y entonces:
	\[
	f(x) - f(x_0) =  \frac{1}{2}(x-x_0)\mathcal{H}f(c)(x-x_0) < 0 \implies f(x) < f(x_0) \forall x \in B(x_0,r)
	\]
\end{enumerate}
\end{proof}


\begin{ndef}[Gradiente]
	Sea $A\subset R^n$ abierto y $f:A\to R$ derivable, llamaremos:
	\[
	\triangledown f(a) = (\frac{\partial f}{\partial x_1}(a),...,\frac{\partial f}{\partial x_n}(a))\in R^n
	\]
	Además, a 
	\[
	\triangledown f: A \to R^n
	\]
	Se le llama campo de vectores.
	\end{ndef}

\textbf{Propeiedad de $\triangledown f$}.

Sea $a \in R$ y $S =\{x \in A :f(x) = a\}$. Entonces, $\forall x_0 \in S \quad \triangledown f(x_0) \  \bot \ S$
\begin{proof}
	Sea $c : I \to R$ con $0 \in I$ una función derivable y $c(I)\subset S$ con $c(0)=x_0$. Ahora,
	\[
	c(t) \in S \quad \forall t \in I \iff f\circ c(t) = f(c(t) = c \implies 0 = (f\circ c)'(t) \implies 0 = Df(c(t))*c'(t)
	\]
	Donde en el último paso hemos usado la regla de la cadena y eso es igual a:
	\[
	Jf(c(t))*c'(t)= \triangledown f(c(t))*c'(t)
	\]
	Y dando a t el valor 0
	\[
	\implies \triangledown f(x_0)*c'(0) = 0
	\]
	Lo que implica la perpendicularidad del gradiente y S.
\end{proof}


\begin{ndef}
	Sea $A\subset R^n$ abierto, no vacío y $a\in A$. Sea también $f\in \mathcal{C}^1(A,\mathbb{R})$. Tomamos las restricciones: $g_1,...,g_k \in \mathcal{C}^1(A,\mathcal{R})$.
	Llamaremos:
	\[
	S=\{ x\in A:g_1(x)=...=g_k(x) = 0\}
	\]
	
\end{ndef}

	

\begin{ndef}[Mínimo relativo condicionado]
	Si $a$ es un punto de $S$, diremos que $f$ presenta un mínimo relativo en $A$ condicionado por las restricciones $g_1(x)=...=g_k(x) = 0$ si y solamente si :
	\[
	\exists \upvarepsilon_0 > 0 : f(a) \leq f(x) \quad \forall x \in A\cap B(a,\upvarepsilon_0) \cap S 
	\]
\end{ndef}



\begin{nth}[de los Multiplicadores de Lagrange]
	En las condiciones que hemos expuesto, las dos definiciones dadas que suponemos ciertas:
	
	Sea $a\in S$ en el que presenta un mínimo relativo condicionado a $S$, entonces $\exists \lambda_0,...,\lambda_k \in R:$
	\begin{nlist}
	\item $(\lambda_0,...,\lambda_k) \ne (0,...,0)$ y $\lambda_0 \geq 0$
	\item  $\lambda_0\frac{\partial f}{\partial x_j}(a)+\sum_{i=1}^k \lambda_i \frac{\partial g_i}{\partial x_j}(a) = 0 \quad \forall j=1,...,N$
\end{nlist}

FALTARLA UN HUECO
\end{nth}
\begin{proof}
	$\exists \upvarepsilon_0 : B(a,\upvarepsilon_0)\subset A $y $f(a) \leq f(x) \quad \forall x \in B(a,\upvarepsilon)\cap S$.
	Vamos a hacer la prueba en dos pasos:
	\begin{enumerate}
	\item $\forall \upvarepsilon \in (0,\upvarepsilon_0) \ \exists M > 0 : f(x) + |x-a|^2 + M \sum_{i=1}^kg_i(x)^2 > f(a) \quad \forall |x-a|= \upvarepsilon$,
	
	Esta prueba se hará por contradicción. Es decir, probando:
	
	$\exists \upvarepsilon > 0 \ \ \forall M> 0 \ \ \exists x\in A: |x-a|=\upvarepsilon$ y $f(x) + |x-a|^2 + M \sum_{i=1}^kg_i(x)^2 \leq f(a) \quad \forall |x-a|= \upvarepsilon$.
	
	Ahora, si tomamos $M=n$, $\exists x_n \in A:$
	\[
	f(x_n) + |x_n-a|^2 + M \sum_{i=1}^kg_i(x_n)^2 \leq f(a) \quad \forall |x-a|= \upvarepsilon
	\]
	Y como $\{x_n\}$ está acotada, existirá $\sigma :N \to N: \{x_{\sigma(n)}\}\to x^{*}$, por lo que 
	\[
	|x^{*}-a|= \lim_{n\to \infty}|x_{\sigma(n)}-a| = \upvarepsilon
	\]
	Tenemos , tras varias transformaciones , que:
	\[
	\frac{f(x_{\sigma(n)})}{-\sigma(n)}+\frac{\upvarepsilon^2}{-\sigma(n)}-\frac{f(a)}{-\sigma(n)} \geq \sum_{i=1}^k g_i(x_{\sigma(n)})^2
	\]
	Pero ahora, los dos segundos términos de la izquierda del $\geq$ convergen a 0 si tomamos límites cuando $n\to \infty$ nos queda que:
	$0 \geq \sum_i g_i (x^*)^2$ y por ello $x^* \in S$
	
	De esta forma, hemos obtenido que $f(x^*) \geq f(a)$ y que
	
	\textbf{FALTAN COSAS QUE NO HE COPIADO}
	
	
	Hemos encontrado por ello $\forall \epsilon \in (0,\epsilon_0)$ un $M>0: f(x) + |x-a|^2 + M \sum_{i=1}^kg_i(x)^2 > f(a) \quad \forall |x-a|= \upvarepsilon$
	\item  $\forall \epsilon \in (0,\epsilon_0) \ \ \exists x^\epsilon$ ( x depende de epsilon, no es x elevado a epsilon) verificando $|x^\epsilon -a| < \epsilon, \ \ \exists(\lambda_0^\epsilon,...,\lambda_k^\epsilon)\in \mathbb{R}^{k+1} : |(\lambda_0^{\epsilon},...,\lambda_k^\epsilon)|= 1$ y $$\lambda_0^\epsilon\{\frac{\partial f}{\partial x_j}(x^\epsilon) + 2(x_j^\epsilon -a_j\}+ \sum_i^k \lambda_i^\epsilon \frac{\partial g_i}{\partial x_j}(x^\epsilon) = 0$$
	
	Utilizando el paso 1, definiendo:
	\[
	F(x):=f(x) + |x-a|^2 + M \sum_{i=1}^kg_i(x)^2
	\]
	Por el teorema de Weierstrass, existe el mínimo de $F(x)$ en $\overline{B}(a,\epsilon)$. Es equivalente a decir:
	\[
	\exists x^\epsilon \in \overline{B}(x,\epsilon) : \ \ F(x^\epsilon)= Min F \ 
	\]
	Y además, por el paso 1:
	\[
	Min F \leq F(A) \leq F|_{\partial B(a,\epsilon)} \implies x^\epsilon \in B(a,\epsilon)
	\]
	Dicho de otra forma, $x^\epsilon$ es un mínimo relativo de F en $B(a,\epsilon)$
	\[
	 \implies \frac{\partial F}{\partial x_j}(x^\epsilon) = 0 \ \ \forall j=1,...,N
	 \]
	 \[
	 \implies \frac{\partial f}{\partial x_j}(x^\epsilon) + 2(x_j^\epsilon -a_j) + M \sum_i^k 2 g_i(x^\epsilon)*\frac{\partial g_i}{\partial x_j}(x^\epsilon) = 0
	 \]
	
	Esa es la derivada de F. Ahora, si dividiéramos esa derivada por:
	\[
	\frac{1}{\sqrt{1+\sum_i^k4M^2g_i(x^\epsilon)^2}}
	\]
	Lo haríamos de módulo 1.
\end{enumerate}
\end{proof}

\textbf{Definición.} Sea $(E,||\cdot||)$ un espacio normado, y sea $G: E \longrightarrow E$. Decimos que $G$ es una \textit{contracción} si es Lipschitziana con constante de Lipschitz menor que 1, es decir, si $\exists c \in (0,1): \ || G(x) - G(y) || \ \le c||x-y|| \ \ \forall x, y \in E$.\\

\textbf{Teorema (del punto fijo de Banach).}
Sea $E$ un espacio de Banach (i.e, un espacio normado y completo), y sea $G:E\longrightarrow E$ una contracción. Entonces, $G$ tiene un único punto fijo, esto es, $\exists !\ x_0 \in E\ tal\ que\ x_0$ = G($x_0$). Además, se verifica lo siguiente: 

\begin{nlist}
\item $\forall y \in E$, si defino inductivamente $y_1 = y$, $y_n = G(y_{n-1})\ \forall n > 1$, entonces $\{y_n\} \rightarrow x_0$. 
\item  $\displaystyle \forall y \in E \ \ ||y - x_o|| \le \frac{||y-G(y)||}{1-c}$
\end{nlist}
\vspace{1em}

\begin{proof} \hfill \\
\underline{Probemos primero $(i)$, y la existencia del punto fijo.} 

Por simplicidad, trabajaremos en un espacio métrico. La primera observación que hacemos es la siguiente: dado $n \in \mathbb{N}$, $n \ge 2$, y usando que $G$ es una contracción, se tiene que
\begin{equation} \label{eq:1} \tag{$\ast$}
d(y_{n+1}, y_{n}) = d(G(y_{n}), G(y_{n-1})) \le c\cdot d(y_{n}, y_{n-1}) \le \dots \le c^{n-1}\cdot d(y_{2}, y_{1})
\end{equation}
Entonces, dados $m,n \in \mathbb{N}$, $m > n$, podemos usar la desigualdad triangular, la desigualdad \eqref{eq:1}, y la fórmula de la suma de los primeros términos de una progresión geométrica, para ver que:
$$0 < d(y_{m}, y_{n}) \leq d(y_m,y_{m-1}) + d(y_{m-1},y_{m-2}) + \cdots + d(y_ {n+1},y_{n})  \leq \left( \sum_{k=n-1}^{m-2}c^k \right) \cdot d(y_2,y_1) = $$ 
\begin{equation} \label{eq:2} \tag{$\ast \ast$}
= \frac{c^{n-1}(c^{m-n} - 1)}{c-1} \cdot d(y_2,y_1) = \left( \frac{c^{n-1}}{1-c} - \frac{c^{m-1}}{1-c} \right) \cdot d(y_2,y_1)
\end{equation}

Como sabemos también que para $c \in (0,1)$ la sucesion $\{\frac{c^{k-1}}{c-1}\} \rightarrow 0$, podemos decir que dicha sucesión es de Cauchy, pues $E$ es completo. Ahora, dado $\epsilon > 0$, escogemos $\epsilon / (d(y_2,y_1) + 1) > 0$, y usando \eqref{eq:2} tenemos que: $$ \exists n_{0} \in \mathbb{N} \ tal\ que\ m,n\geq n_{0} \implies d(y_{m}, y_{n}) \leq \left| \frac{c^{n-1}}{1-c} - \frac{c^{m-1}}{1-c} \right| \cdot d(y_2,y_1) < \frac{\epsilon}{d(y_2,y_1) + 1} \cdot d(y_2,y_1) < \epsilon $$ 

Esto nos garantiza que $\{y_{n}\}$ es de Cauchy. Entonces, como el espacio $E$ es completo, tenemos que $\{y_n\}$ es convergente, es decir, $\exists x_0 \in E:\ \{y_{n}\}\rightarrow x_0$. Veamos ahora que $x_0$ es necesariamente un punto fijo de $G$.

En efecto, por ser $x_0$ el límite de la sucesión, dado $\epsilon > 0, \exists n_0 \in \mathbb{N}: n \ge n_0 \Rightarrow d(y_n, x_0) =$ $= d(G(y_{n-1}),x_0) <  \epsilon$, lo que implica que $\{G(y_{n})\}\rightarrow x_0$. Pero $\{G(y_{n})\}\rightarrow G(x_0)$ por ser G continua (es Lipschitziana). Por tanto, $x_0 = G(x_0)$ por unicidad del límite.
\vspace{0.5em}

\underline{Veamos ahora la unicidad del punto fijo.}

\vspace{0.5em}
Supongamos que existe otro punto fijo $y_0 \in E$. Entonces, como $G$ es una contracción, se tendría que: 
$$d(G(y_0),G(x_{0})) \le c\cdot d(y_0, x_0) \Rightarrow d(y_0, x_0) \le c \cdot d(y_0, x_0)\ con\ 0 < c < 1$$ Esto implica que $d(y_0,x_0) = 0 \Rightarrow x_0 = y_0$. Por tanto, si hay dos puntos fijos, han de ser necesariamente iguales.\\

\underline{Por último, probemos $(ii)$.}

Sea $y \in E$. Entonces, si tomamos $y_1$ = $y$ en la sucesión, se tiene que, usando \eqref{eq:2}: $$d(y,x_0) = \lim_{n\rightarrow \infty}d(y_1,y_n) \le \lim_{n\rightarrow \infty}\left(\sum_{k=0}^{n-2} c^k\right)\cdot d(y_1, y_2)= \lim_{n\rightarrow \infty}\left(\sum_{k=0}^{n-2} c^k\right)\cdot d(y, G(y)) = \frac{d(y, G(y))}{1-c} $$

donde hemos usado la fórmula de la suma de los términos de una \textit{serie geométrica} de razón $0 < c < 1$.
\end{proof}

\section{Teorema de la función inversa e implícita}
\begin{nth}[Teorema de la función inversa. Caso N=1]
	Sea $A\subset R$, $a \in A$ abierto y $f:A \to \R$. Supongamos que $f$ es de clase 1 y $f'(a) \ne 0$.\\
	 Entonces, $\exists I $ un intervalo abierto con $I \subset A$ y $a \in I$, y $\exists J$ intervalo abierto con $J\subset \R$ con $f(a) \in J$ tal que $f|_I: I \to J$ es invertible si $g = (f|_U)^{-1} \implies g'(f(a)) = \frac{1}{f'(a)}$
\end{nth}

Ahora, como podemos ver una relación: 
\[
Df(a) \equiv Jf(a) = ( f'(a))
\]
Por lo que, el determinante de $Jf(a) \ne 0$, es decir que $Df(a)$ es invertible
\begin{nth}[Teorema de la función inversa. Caso general.]
	Sea $V \subset \R^n$ abierto y $a\in A$ y $f:V \to \R^n$ de clase 1. Si  
$Df(a): \R^n \to \R^n$ es invertible \\
$\implies \exists U\subset V$ abierto y $\exists W\subset\R^n$ abierto tales que $a\in U $, $f(a) \in W$ y $f|_U : U \to W$ es biyectiva, y por tanto existe la inversa de la función $g = (f|_U)^{-1}: W \to U$ es de clase 1 y:
	\[
	Dg(f(a)) = (Df(a))^{-1}
	\]
Recordemos que decir que $Df(a)$ es invertible es lo mismo que decir que $det(Jf(a)) \ne 0)$
Y decir que $Dg(f(a)) = (Df(a))^{-1}$ es lo mismo que decir $Jg(f(a)) = (Jf(a))^{-1}$
\end{nth}
\begin{proof}
	El caso complicado de la demostración estaría en el caso en el que $0 = a$ y $f(a) = 0$. Ahora, si pudiéramos demostrar que la derivada es la aplicación identidad, entonces podremos demostrarlo en general. Lo que vamos  a probar es lo siguiente:
	
	\begin{nth}[Caso particular]
		Sea $V \subset \R^n$ abierto y $0=a\in A$ ,y $f:V \to \R^n$ de clase 1 con $f(0) = 0$. Si  
$Df(a): \R^n \to \R^n$ es invertible \\
$\implies \exists U\subset V$ abierto y $\exists W\subset\R^n$ abierto tales que $0\in U $, $0 \in W$ y $f|_U : U \to W$ es biyectiva, y por tanto existe la inversa de la función $g = (f|_U)^{-1}: W \to U$ es de clase 1 y:
	\[
	Dg(0) = I
	\]
\end{nth}
\begin{proof}[Demostración del caso particular del teorema]\hfill \\
Sea V un abierto de $\R^n$ con $0 \in V$. Como V es abierto, $2\delta > 0$ tal que $B(0,2\delta) \subset V$. Además, como f es de clase 1, su derivada es continua.
Ahora, como $Df(0)=I$, $\forall \epsilon > 0 \ \exists \delta > 0$ tal que $B(0,2\delta) \subseteq V$ y $||Df(x)-I|| < \epsilon \ \forall x \in B(0,2\delta)$. Hay que recalcar que $\delta$ depende de $\epsilon$.Vamos a suponer también que para epsilon más pequeños, el delta será también más pequeño. De hecho, vamos a escribir $\delta(\epsilon)$

\begin{nota}
	[1]
	Tomando ahora $r= 2\delta(\frac{1}{2}) \implies f$ es inyectiva en $B(0,r)$
\end{nota}
\begin{proof}[Demostración de la afirmación 1]\hfill \\
Sean $x,y \in B(0,r)$ tales que $f(x) = f(y)$. Entonces, por el TVM, $\exists z \in B(0,r)$ tal que $f(x)-f(y) = Df(z)(x-y) \implies Df(z)(x-y) = 0$.\\
Es fácil ver que $||Df(z) -I || < \frac{1}{2}$ pues $z \in B(0,r) = B(0,2\delta(\frac{1}{2})$\\
	Ahora, si tomamos\\
	$ 0=(y-x)*Df(z)(y-x) = (y-x)[Df(z)(y-x) -I(y-x) +(y-x)] = (y-x)[(Df(z)-I)(y-x) +(y-x)] = ||y-x||^2 +(y-x)(Df(z)-I)(y-x)$\\ 
	Y , como es una aplicación lineal, usamos que es Lipschitziana y nos queda que eso es mayor o igual que
y por la desigualdad de Cauchy Schwarz\\
	$ ||y-x||^2 +(y-x)(Df(z)-I)(y-x) \geq || y-x||^2 - ||Df(z) -I|| \ || y-x||^2 \implies ||y-x|| = 0$ por lo que $x=y$, por lo que $f$ es inyectiva.


\end{proof} %de la nota 1
Ahora, tenemos que probar que es sobreyectiva, habría ver si $\forall h \in W$ $\ \exists x \in U =B(0,r) : \ \ f(x) = h$.
Para empezar, podemos escribir eso como $x+f(x)-x = h \implies x = h+(x-f(x) \implies x = h +r(x)$ siendo $r(x)=(x-f(x))$ y luego, llamaremos $G(x) = h+r(x)$, quedándonos $x=G(x)$. Así, $G: U \to \R^n$ con U abierto de $\R^n$ y así lo que tendríamos que encontrar es un $x\in U$ tal que $G(x) = x$, y si existe a ese $x$ lo llamaremos punto fijo de G.

Ahora, por el teorema del punto fijo de Banach o teorema de la contracción

\textbf{Lema}.
$\forall \epsilon \in (0,1/2)$ se tiene:
\begin{nlist}
	\item $f(\overline{B(0,\delta(\epsilon)}) \supseteq  B(0,\delta(\epsilon)/2)$
	\item \[\begin{rcases}
	h \in B(0, \delta(\epsilon)/2)\\
	x \in \overline{B(0, \delta(\epsilon))} :  \ f(x) = h
\end{rcases} \implies ||x-h|| \leq \frac{\epsilon}{1-\epsilon}||h||\]
\end{nlist}

\begin{proof}[Demostración del lema]
	\begin{nlist}
	\item Sea $h \in B(0, \delta(\epsilon)/2) $. Definimos $G(y) =  y-f(y) +h$, con $G: \overline{B(0, \delta(\epsilon)/2)} \to \R^n$.
	
	Si $x \text{ fijo } \in G \iff x = G(x) = x-f(x) +h \iff f(x) = h$. Definimos también $r(y):= y - f(y)$.
	
	Ahora, vamos entonces a probar que $G$ es una contracción para poder usar el teorema del punto fijo de Banach. Tenemos que:
	\[
	DG(y) = I - Df(y) = Dr(y)
	\]
	Y tenemos también que $||DG(y))|| \leq \epsilon \quad \forall y \in \overline{B(0, \delta(\epsilon)/2)}$, luego si hacemos $G(y) - G(\overline{y})$, eso es:
	\[
	||G(y) - G(\overline{y})|| \leq ||DG(z)||\ ||y-\bar{y}|| \leq \epsilon ||y-\bar{y}||
	\]
	Con $y,\bar{y} \in \overline{B(0, \delta(\epsilon)/2)} \implies  $ G es una contracción en $\overline{B(0, \delta(\epsilon)/2)}$
	Y como el teorema del punto fijo de Banach dice que la función G va de un espacio métrico en el mismo, tenemos que probar que $G(\overline{B(0, \delta(\epsilon)/2)}) \subset \overline{B(0, \delta(\epsilon)/2)}$.
	
	Probamos esto, vemos que $G(0) = 0-f(0) + h = h$, y ahora $\forall x \in \overline{B(0, \delta(\epsilon)/2)} \quad ||G(x)|| \leq ||G(x)-G(0)|| +||G(0)||$, y como es una contraccion, tenemos que eso es $\leq ||x|| + || h || < \epsilon \delta(\epsilon) + \frac{\delta(\epsilon)}{2} < \delta (\epsilon)$, lo que significa que $G(x) \in B(0, \delta(\epsilon)/2)$
	
	Ahora, sí podemos usar el teorema del punto fijo de Banach, $\exists ! x \in \overline{B(0, \delta(\epsilon)/2)}: \ \ G(x) = x \implies f(x) = h$, como queriamos probar
	\item Deducimos que $||x-h || \leq \frac{||h-G(x)||}{1-\epsilon} \frac{1}{1-\epsilon} || h - G(h)|| = \frac{1}{1-\epsilon} || -r(y)|| =  \frac{1}{1-\epsilon} || r(h)|| \frac{1}{1-\epsilon} || r(h) - r(0)|| \leq \frac{1}{1-\epsilon} \epsilon|| h|| = \frac{1}{1-\epsilon} ||h||$ 
	
	Y queda así probado el lema
\end{nlist}
\end{proof} %del lema

Ahora, para probar que la f es biyectiva que era lo que queríamos probar del teorema de la función inversa, tenemos que fijar un epsilon que será cualquier número entre $0 $ y $1/2$, pues hemos obtenido un lema para ese rango de $\epsilon$.

Consideramos el conjunto: $W = B(0, \delta(\epsilon)/2)$


Podemos ver que $\delta(\epsilon) < \delta(1/2) < 2 \delta(1/2) = r$. Ahora, usando el primer apartado del lema obtenemos que:
\[
f(\overline{B(0,\delta(\epsilon)}) \supseteq B(0, \delta(\epsilon)/2)
\]
Y consideramos:
\[
U= B(0, \delta(\epsilon)) \cap f^{-1}(W)
\]
y, como f era inyectiva en la $B(0,r)$, si definimos: $f: U \to W$ es inyectiva porque es una restricción de la f inicial y es sobreyectiva por la definición de U.

Así, hemos probado que $f|_U : U \to W$ es biyectiva, lo segundo que queríamos probar.

Nos queda por último comprobar la derivabilidad de la función inversa. Habría que probar que:
\[
\lim_{h\to 0}\frac{||g(h) - g(0) -I(h) ||}{||h|| } = 0 \iff \lim_{h \to 0}\frac{||g(h) - h||}{||h||}
\]


\end{proof} %del caso particular
\end{proof} %del caso general









El teorema de la función inversa lo usaremos si tenemos las mismas ecuaciones que incógnitas. Si tenemos menos ecuaciones que incógnitas, usaremos el teorema de la función implícita.

\begin{nth}[Teorema de la función implícita]
	Sea $A \subset \R^n x\R^m$ abierto y no vacío. Sea $F: A \to \R^n$ una función de clase 1. Fijando un $(x_0,y_0) \in A$ tal que $F(x_0,y_0) = 0$. Si:
	\[
	det(\begin{pmatrix}
 \frac{\partial F_1}{\partial n_1}(x_0,y_0) & \cdots & \frac{\partial F_1}{\partial x_n}(x_0,y_0) \\
 \cdots& & \cdots\\
 
 \frac{\partial F_n}{\partial x_1}(x_0,y_0) & \cdots & \frac{\partial F_n}{\partial x_n} (x_0,y_0)
\end{pmatrix}) \ne 0
	\]
	
	Entonces, existe $U$ un entorno abierto de $x_0$ en $\R^n$ y también existe $V$ un entorno abierto de $y_0$ en $\R^m$ tal que $UxV \subset A$ y $\forall y \in V \ \ \exists! x \in U : F(x,y) = 0$
\end{nth}
\begin{proof}
	Lo demostraremos a partir del teorema de la función inversa.\\
	Sea $f:\R^nx\R^m \to \R^nx\R^m$ definida por $f(x,y) = (F(x,y),y-y_0)$.
\end{proof}
\end{document}